% 제2장 이론적 배경
% 1차 심사 반영: "이론적 배경 및 연구 프레임워크" → "이론적 배경" (축약)

\chapter{\vrev{이론적 배경}}
\label{ch:theory}

\vrev{%
\chg{1.2(연구의 필요성)에서 제기한 정적 RTO의 세 한계와 통합적 접근의 필요성에 따라서,
제2장(이론적 배경)은 동적 복구목표시간($\DRTO$) 모델과 그 조건 체계, 연구 가설을 선행
이론에 근거하여 처음으로 도출하고 확립한다.}
본 장에서는 BCMS와 RTO의 기존 한계,
관측성 이론, 불확실성 모델링, 시그모이드 함수와 바운딩 기법에 대한
선행연구를 검토하고, 본 연구의 이론적 토대를 정립한다.
이어서 관측성 \m{기반 $\DRTO$ 모델} 구조와 수학적 성질을 제시하고,
\jrev{10개 \m{연구 가설}}을 체계적으로 수립한다.
}

\section{BCM(업무연속성관리)}
\label{sec:lit-bcm}
%% ===================================================================

\chg{본 연구가 다루는 RTO가 자리하는 제도적 맥락을 분명히 하기 위해, BCM의 개념과 역사적 발전, 국제 표준 체계, 그리고 RTO가 도출되는 BIA 절차를 차례로 검토한다.}

\chg{본 연구에서 사용하는 BCM과 BCMS의 의미를 구분하여 밝혀 둔다. \m{BCM은} 조직이 재해와 장애, 위협으로부터 핵심 업무를
지속하고 복구하기 위한 관리 활동과 개념 전반을 가리킨다. 이에 비해
\m{BCMS는} 이러한 BCM 활동을
정책과 절차, 자원, 문서, 평가의 형태로 체계화한 경영 시스템을 가리키며,
ISO 22301:2019 인증의 직접 대상이 된다. 따라서 본 연구는 일반적 활동이나
개념을 가리킬 때는 BCM을, 정형화된 시스템이나 표준, 인증을 가리킬 때는 BCMS를
사용하여 일관되게 구분한다.}
\citet{herbane2010disaster}\m{은} BCM의 역사적 발전 과정을 세 단계로 구분하였다.
첫째, 1970--1980년대의 IT 재해복구(Disaster Recovery) 중심 단계에서는 데이터 백업과
시스템 이중화 등 기술적 복구에 초점을 두었다.
둘째, 1990년대에는 \m{BIA의 도입}과 함께
IT를 넘어 전사적 업무 프로세스로 관리 범위가 확대되었다.
셋째, 2000년대 이후에는 \citet{bhamra2011resilience}이 정리한
\textbf{조직 복원력}(organizational resilience) 관점에서
BCM을 전략적 경영 기능으로 통합하는 방향으로 진화하였다.

이러한 발전의 제도적 결실로서 영국 BSI는 2006년에 BS~25999를 발표하여
BCM의 국가 표준 단계의 핵심 프레임워크를 제시하였다\citep{bs25999}.
이후 국제표준화기구(ISO)는 2012년에 \nrev{ISO 22301:2012}를 제정하였으며,
2019년에 개정판을 공표하여 현재의 국제 표준으로 자리매김하였다.
\nrev{ISO 22301:2019}는 BCMS의 수립·구현·유지·지속적 개선을 위한 요구사항을 \m{규정하며}, Plan--Do--Check--Act(PDCA) 주기에
기반한 관리체계를 \m{요구하고 있다}.
\nrev{ISO 22313:2020}은 이에 대한 실행 지침서로서 BIA 수행 방법과
복구 전략 수립 절차를 구체적으로 \m{안내한다}.

\figref{fig:bcm-pdca}\refpo{fig:bcm-pdca}{은}{는} \nrev{ISO 22301:2019}의 PDCA 주기에서
RTO가 결정되는 위치를 도식화한 것이다. 정적 RTO는 Do(운영) 단계의 BIA에서 한 번
결정되어 다음 BIA까지 고정되는 반면, 본 연구의 $\DRTO$는 운영과 모니터링이
이루어지는 Do--Check 구간에서 실시간으로 파악된다.

\begin{figure}[htbp]
  \centering
  \begin{tikzpicture}[
    phase/.style={draw, semithick, rounded corners=8pt, minimum width=34mm, minimum height=16mm,
                  align=center, font=\small, inner sep=4pt},
    cyc/.style={-{Stealth[length=3.2mm]}, very thick, gray!55!black},
    annot/.style={draw=red!55, fill=red!7, semithick, rounded corners=4pt,
                  align=center, font=\footnotesize, inner sep=5pt},
    darr/.style={-{Stealth[length=2.8mm]}, very thick, dashed, red!70!black}
  ]
    % --- PDCA 순환 (시계방향: 상→우→하→좌) ---
    \node[phase, fill=blue!10]   (plan)  at (0,  2.6) {\textbf{Plan}\\[1pt]\footnotesize BCMS 수립};
    \node[phase, fill=green!13]  (do)    at (4.2, 0) {\textbf{Do}\\[1pt]\footnotesize BIA 수행\\\footnotesize \textbf{정적 RTO 결정}};
    \node[phase, fill=yellow!22] (check) at (0, -2.6) {\textbf{Check}\\[1pt]\footnotesize 모니터링, 측정};
    \node[phase, fill=orange!14] (act)   at (-4.2, 0) {\textbf{Act}\\[1pt]\footnotesize 지속적 개선};
    \node[font=\small\bfseries, align=center] at (0,0) {ISO 22301:2019\\PDCA 주기};
    \draw[cyc] (plan.east)  to[bend left=22] (do.north);
    \draw[cyc] (do.south)   to[bend left=22] coordinate[midway] (dc) (check.east);
    \draw[cyc] (check.west) to[bend left=22] (act.south);
    \draw[cyc] (act.north)  to[bend left=22] (plan.west);
    % --- 정적 RTO vs 본 연구의 DRTO (우측 분리 배치, 직각 연결) ---
    \node[annot] (static) at (9.3, 2.6)
      {\textbf{정적 RTO}\\BIA에서 결정};
    \draw[darr] (static.west) -| (do.north);
    \node[annot, fill=red!11] (drto) at (9.3, -2.6)
      {\textbf{본 연구: $\DRTO$}\\운영--모니터링 중\\실시간 동적 파악};
    \draw[darr] (drto.west) -| (dc);
  \end{tikzpicture}
  \caption{\nrev{ISO 22301:2019} PDCA 주기 기반 RTO 결정 위치}
  \label{fig:bcm-pdca}
\end{figure}


따라서 BIA는 BCMS의 핵심 구성 요소로서 중단 사태가 조직에 미치는 영향을 분석하고,
핵심 업무 기능의 복구 우선순위와 복구 목표 시간을 결정하는 기초를 제공한다
\citep{cerullo2004bcp, snedaker2013bcdr}.
\citet{gibb2006bcm}은 BIA 결과가 조직의 복구 전략과 자원 배분에 직접적으로 연결되어야
함을 \m{강조하며}, 정보관리 관점에서의 BCM 프레임워크를 제안하였다.

RTO는 중단 사태 발생 이후 핵심 업무 기능을
사전에 정의된 수준까지 복구하는 데 소요되는 목표 시간을 의미한다.
\nrev{ISO 22301:2019}에서는 RTO를 ``제품이나 서비스의 납품 또는 수행, 활동의 재개,
또는 자원의 복구가 이루어져야 하는 시점까지의 기간''으로 정의한다.
미국 NIST SP~800-34~Rev.~1에서는 RTO를 ``시스템 중단 이후 시스템과 관련 자원이
복구되어야 하는 최대 허용 시간''으로 \m{규정하고 있다}\citep{nist800-34r1}.

MTPD는 조직이 중단 사태를
견딜 수 있는 최대 허용 기간을 의미하며, RTO는 반드시 MTPD 이내에 설정되어야 한다.
즉, $\text{RTO} \leq \text{MTPD}$의 관계가 성립한다\nrev{(ISO 22313, 2020; \citeauthor{snedaker2013bcdr}, \citeyear{snedaker2013bcdr})}.
\citet{hiles2010bcm}은 BIA에서 도출된 MTPD가 RTO 설정의 상한 기준임을
강조하였다.

국내 BCM 연구는 \chg{BIA, 위험평가, 평가지표} 측면에서 점진적 진전을 이루었으나, RTO의 동적 설정 문제는 다루지 않았다.
\citet{cheung2019quickhit}는 BIA 및 위험평가를
효율적으로 수행하기 위한 Quick-Hit 프레임워크를 제안하여 BCMS 도입 초기의
실무적 장벽을 낮추고자 하였다. \citet{jang2021bcms}은 BCMS의 위험평가가
경영성과에 미치는 영향을 실증 분석하여 BCM과 조직 성과 간의
정(+)의 관계를 확인하였으며, \citet{jung2023bcms}은 BCMS의 체계적 평가를 위한
평가지표를 개발하여 국내 조직의 BCM 성숙도 진단에 기여하였다.

\chg{이처럼 BCM 연구는 체계 도입과 성과 평가로 꾸준히 확장되어 왔으나, 그 핵심 통제 변수인 RTO를 어떻게 설정할 것인가의 문제는 여전히 정적 관점에 머물러 있다. 2.2(RTO의 정적 한계)에서는 이 정적 RTO 설정이 지닌 구조적 한계를 살펴본다.}
  % 2.1 업무연속성관리(BCM)


%% ===================================================================
\section{RTO의 정적 한계}
\label{sec:lit-rto-static}
%% ===================================================================

\chg{앞 절에서 살펴본 BCM의 핵심 산출물 가운데 하나가 RTO이다. 그러나 현행 RTO는 복구 환경의 변화와 무관하게 단일 정적 값으로 고정된다는 구조적 제약을 안고 있다. 이 절에서는 그 한계를 정적 설정의 구조적 문제와 선행연구의 접근을 중심으로 검토한다.}

%% -------------------------------------------------------------------
\subsection{정적 RTO 설정의 구조적 문제}
\label{subsec:static-rto-problem}
%% -------------------------------------------------------------------

전통적 RTO 설정 방식에는 세 가지 근본적 한계가 존재하며, 이는 본 논문
\secref{sec:problem_statement}(연구 필요성)에서 제시한 세 가지 연구 필요성과 직접 대응된다.

\chg{첫째는 관측성 미반영이다. 정적 RTO는 BIA에서 도출된 고정값으로서 조직의 운영
가시성, 곧 관측성(observability)을 반영하지 못한다. 시스템 복잡도와 인력 역량, 기술 환경 등
운영 환경이 시간과 상황에 따라 변화하는 현실에서 관측성이 높은 시스템은 장애 원인을 신속히
파악하여 복구 시간을 단축할 수 있으나, 전통적 BIA는 이러한 운영 역량 차이를 RTO에
반영하는 메커니즘을 갖추지 못한다. 그 결과 관측성이 높은 조건에서는 복구 자원의 과다
예비(over-provisioning)가, 낮은 조건에서는 과소 예비(under-provisioning)가 동시에 초래된다.
둘째는 일상적 불확실성의 정량화 부재이다. 복구 과정에서 발생하는 장비 가용성과 인력 동원 속도,
외부 공급망 의존도 같은 일상적 불확실성이 정량적으로 모델링되지 않는다. 실제 복구 시간은 다양한
예측 불가능한 요인에 의해 목표 시간을 초과할 수 있으나, 기존 접근법은 이러한 정규분포(Gaussian)적
변동을 체계적으로 반영하지 못한다. 셋째는 극단 사건의 미모형화이다. 사이버공격에 의한 데이터
센터 전면 장애나 복합 재난으로 인한 공급망 동시 중단처럼 일상적 변동을 넘는 극단적 사건에 의한
복구 지연은, 정규분포 가정으로는 포착할 수 없는 파레토(Pareto, 중꼬리, heavy-tail) 분포 특성을
보인다~\citep{coles2001extremes, embrechts1997extremal, taleb2007blackswan}.}


%% -------------------------------------------------------------------
\subsection{선행 연구에서의 RTO 접근}
\label{subsec:prior-rto-approaches}
%% -------------------------------------------------------------------

기존의 RTO 관련 연구는 대체로 BIA 방법론의 개선, RTO 달성을 위한 기술적 수단
(이중화, 백업 전략), 또는 특정 산업 분야에서의 사례 분석에 집중하고 있다.
\citet{snedaker2013bcdr}는 IT 재해복구와 업무연속성 계획 수립의 체계적 방법론을
제시하였으나, RTO 자체를 동적으로 조정하는 수리 모델은 다루지 않았다.
\citet{cerullo2004bcp}는 금융 서비스 산업에서의 BCP 수립 사례를 분석하여
RTO의 실무적 중요성을 강조하였으나, RTO의 동적 특성에 대한 논의는 부재하였다.

\chg{최근에는 업무연속성과 재해복구를 통합하여 조직 복원력 관점에서 정량적으로 접근하려는 시도가 이어지고 있다. \citet{sahebjamnia2015integrated}는 핵심 업무의 재개와 복구를 동시에 고려하는 다목적 혼합정수계획 모형으로 복구 자원의 최적 배분 문제를 정식화하였고, \citet{torabi2016enhanced}는 최선-최악 방법(Best--Worst Method)에 기반한 향상된 위험평가 체계를 제안하여 BCM의 위험 정량화 수준을 끌어올렸다. 그러나 이들 정량적 프레임워크에서도 RTO 자체는 위험평가와 자원배분의 고정된 입력 매개변수로 취급되며, 운영 환경의 관측성이나 복구 시간의 확률적 변동에 따라 RTO를 동적으로 조정하는 메커니즘은 여전히 다루어지지 않는다.}

\citet{cheung2019quickhit}, \citet{jang2021bcms} 그리고 \citet{jung2023bcms}을 포함한 국내 BCM
연구 또한 표준 적용과 조직 성과 평가에 집중하고 있어, RTO의 동적 조정이나
운영 조건에 따른 적응적 복구 모델에 대한 연구가 상대적으로 부재함을 드러낸다.
\citet{lee2026drto}이 관측성 채널(C0$\to$C1)을 도입하여
초기 접근을 수행하였으나, 불확실성 채널(C2/C3)의 체계적 통합은 향후 과제로
남아있었다.
이러한 세 가지 한계와 향후과제가 본 연구에서 $\DRTO$ 프레임워크를 통해 해소하고자 하는
핵심 문제의식의 출발점이다.

\chg{이러한 적응적 복구를 실현하려면 복구 환경의 상태를 실시간으로 관측할 수 있어야 하며, 다음 절에서는 그 이론적 토대인 관측성 개념을 검토한다.}
  % 2.2 복구목표시간(RTO)의 정적 한계


%% ===================================================================
\section{관측성(Observability) 이론}
\label{sec:lit-observability}
%% ===================================================================

\chg{정적 RTO의 한계를 극복하려면 복구 환경의 상태 변화를 실시간으로 파악할 수 있어야 한다. 이러한 상태 파악 능력을 이론적으로 뒷받침하는 개념이 관측성(observability)이며, 이 절에서는 제어이론에서 출발한 관측성 개념과 그 현대적 확장을 검토한다.}

%% -------------------------------------------------------------------
\subsection{제어 이론에서의 관측성 (Kalman, 1960)}
\label{subsec:kalman-observability}
%% -------------------------------------------------------------------

관측성(observability)의 개념은 \citet{kalman1960control}이 제어 이론에서 체계적으로 형식화하였다.
Kalman은 선형 시불변(Linear Time-Invariant, LTI) 시스템
\begin{equation}
  \dot{x}(t) = Ax(t) + Bu(t), \quad y(t) = Cx(t)
  \label{eq:lti-system}
\end{equation}
에서 시스템의 출력 $y(t)$를 관찰함으로써 초기 상태 $x(0)$를 유일하게 결정할 수 있는지를
판별하는 기준으로 관측성 행렬(observability matrix)의 계수(rank) 조건을 제시하였다.
관측성 행렬
\begin{equation}
  \mathcal{O} = \begin{bmatrix} C \\ CA \\ CA^2 \\ \vdots \\ CA^{n-1} \end{bmatrix}
  \label{eq:obs-matrix}
\end{equation}
이 완전 계수(full rank)일 때 시스템이 완전 관측 가능(completely observable)하다.

이 개념은 시스템의 내부 상태를 외부 출력만으로 추론할 수 있는 역량을 수학적으로
정의한 것으로, 단순한 모니터링(monitoring)과는 본질적으로 구별된다~\citep{beyer2016sre}.
모니터링이 사전에 정의된 지표의 임계값 초과 여부를 감시하는 반응적(reactive) 방식인 반면,
관측성은 시스템 출력의 충분성(sufficiency)에 관한 구조적 속성이다.

\chg{Kalman 이후 제어 이론에서 관측성은 상태 추정(state estimation)의 토대로 발전하였다. \citet{luenberger1966observers}는 관측성 조건이 충족되는 선형 시스템에서 직접 측정되지 않는 내부 상태를, 가용한 출력과 입력만으로 구동되는 보조 동역학계인 상태 관측기(state observer)를 통해 재구성할 수 있음을 보였다. 이로써 관측성은 가부 판별을 넘어 관측 가능한 시스템에서는 내부 상태를 실제로 복원할 수 있다는 구성적 의미를 획득하였고, 이후 상태 추정과 피드백 제어의 핵심 원리로 자리 잡았다. 본 연구가 관측성을 복구 환경의 내부 상태에 대한 가시성으로 해석하는 것은, 바로 이 출력으로부터 내부 상태를 복원하는 역량이라는 제어 이론적 함의를 BCM 영역으로 확장한 것이다.}


%% -------------------------------------------------------------------
\subsection{BCM 영역으로의 확장}
\label{subsec:observability-bcm}
%% -------------------------------------------------------------------

%% --- BCM 맥락의 정규화 ---
\subsubsection{BCM 맥락에서의 관측성 정규화}
\label{subsubsec:obs-normalization}

본 연구에서는 Kalman의 제어 이론적 관측성 개념을 BCM 영역으로 확장하여 적용한다.
\chg{Kalman의 원래 정식화는 관측성 행렬의 계수 조건을 통해 시스템이 ``완전 관측
가능한가''를 가능과 불가능의 이분법으로 판별한다. 그러나 \chg{BCM과 재해복구} 실무에서 조직의 관측성은 이분법으로 주어지지 않는다.
\chg{현대 IT 운영에서 관측성은 로그(log), 메트릭(metric), 트레이스(trace)라는 세 가지
텔레메트리(telemetry)를 통해 확보되며~\citep{majors2022observability}, 조직마다 센서와
모니터링 도구의 구비 정도, 그리고 이 세 텔레메트리의 수집과 보존 수준이 달라 그 확보
정도가 연속적으로 변한다.} 그 결과 조직의 관측성은 0과 1 사이에서 연속적으로 분포한다. 본 연구는 바로 이 부분을 확장하여 적용하는데, Kalman의
이진 판별 대신 관측성의 \emph{정도}를 $O_{\text{raw}} \in [0, 1]$의 연속 점수로 정량화한다.
이렇게 연속화하는 이유는, 복구 시간의 단축 효과가 관측성의 유무가 아니라 그
수준에 연속적으로 의존한다고 보기 때문이며, 연속 점수라야 시그모이드 바운딩을 통한 복구 시간
단축량의 정량적 모형화가 가능해지기 때문이다.}
이때 $O_{\text{raw}} = 0$은 시스템 상태를 전혀 파악할 수 없는 상태를,
$O_{\text{raw}} = 1$은 완전히 관측 가능한 상태를 의미한다.


%% --- 관측성-복구 능력 구조적 연결 ---
\subsubsection{관측성과 복구 능력의 구조적 연결}
\label{subsubsec:obs-recovery-link}

\citet{avizienis2004basic}와 \citet{hollnagel2006resilience}이 각각
의존성과 복원공학 문헌에서 논의한 바와 같이, 관측성과 복구 능력 사이에는
구조적 연결이 존재한다.
시스템 상태를 보다 정확히 파악할 수 있을수록 장애의 근본 원인(root cause)을
신속히 진단할 수 있고, 이에 따라 복구에 소요되는 시간이 단축된다.
\citet{avizienis2004basic}는 가용성(availability), 신뢰성(reliability),
안전성(safety), 무결성(integrity)을 포괄하는 의존성(dependability)의 분류 체계에서
장애 진단 능력이 시스템의 가용성에 직접적으로 영향을 미침을 논의하였다.
\citet{hollnagel2006resilience}는 복원공학(Resilience Engineering) 관점에서 시스템의
예측, 감시, 대응, 학습 능력을 강조하며, 이 중 감시(monitoring) 능력이 본질적으로
관측성에 해당한다고 볼 수 있다.
\chg{본 연구는 이 두 이론을 관측성과 복구 시간 사이의 단조 감소 관계로 반영하여 적용한다.
\citet{avizienis2004basic}에서는 장애 진단 능력이 가용성에 기여한다는 인과를,
\citet{hollnagel2006resilience}에서는 감시 능력이 복원력의 핵심이라는 명제를 각각
가져온다. 본 모델은 이 ``관측성에서 진단으로, 진단에서 복구 시간 단축으로''의 인과를
관측성 채널의 하향 조정량 $E_{O,\text{bnd}} \leq 0$으로 구현한다. 즉 관측성이 높을수록
$\DRTO$를 기준 RTO $R_0$ 아래로 끌어내리는 음의 효과로 정량화하는 것이다. 두 이론을
이렇게 적용한 이유는, 이론이 제시한 인과 관계를 모델에서 관측성과 복구 시간 사이의
단조 감소 관계로 직접 반영함으로써, 관측성 투자에 따른 복구 시간 단축을 정량적으로
예측할 수 있게 하기 위함이다.}


%% --- SRE와 세 가지 핵심 요소 ---
\subsubsection{\chg{현대 IT 운영의 관측성 SRE와 세 가지 핵심 요소(three pillars)}}
\label{subsubsec:obs-sre-pillars}

현대 IT 운영에서의 관측성은 Kalman의 제어 이론적 개념을 분산 시스템
환경으로 재해석한 것으로 볼 수 있으며, \citet{beyer2016sre}와
\citet{majors2022observability}이 이를 SRE 실무 체계로 정착시켰다.
\citet{beyer2016sre}는 Google의 SRE(Site Reliability Engineering)
실무를 통해 모니터링과 관측성의 차이를 구분하였고,
\citet{majors2022observability}는 관측성의 세 가지 핵심 구성 요소(three pillars), \chg{곧
메트릭(metrics)과 로그(logs), 트레이스(traces)를} 체계화하였다.
이 세 구성 요소의 통합적 활용을 통해 조직은 예상치 못한 장애 상황에서도
시스템의 내부 상태를 효과적으로 파악할 수 있다.
\p{[}그림~\ref{fig:obs-three-pillars}\p{]}는 이러한 세 가지 핵심 요소(three pillars)의 구성 관계를 도식화한 것이다.


%% --- 운영화 매핑 ---
\subsubsection{\chg{데이터 유형에서} \chg{도구와 목적} \chg{차원으로의 운영화 매핑}}
\label{subsubsec:obs-operationalization}

이러한 세 가지 핵심 요소(three pillars) 분류는 관측성을 구성하는 \textbf{데이터 유형(data primitives)} 차원의
정의이다. 본 연구는 \chg{BCM과 재해복구} 맥락의 운영화 단계에서 관측성을 \chg{도구와 목적}별 통합 시스템
차원, \chg{곧 APM(Application Performance Monitoring, 서비스 신뢰성 영역)과 SIEM(Security
Information and Event Management, 보안 이벤트 대응 영역)의 두 축으로} 재구성하여 정량화한다.
APM과 SIEM은 각각 \chg{메트릭, 로그, 트레이스}를 통합 활용하되 측정 목적과 주된 데이터 출처가
상이하므로 분리 산출 후 가중 합성($O_{\text{raw}} = \alpha \cdot O_{\text{APM}}
+ (1-\alpha) \cdot O_{\text{SIEM}}$)으로 통합하며, 이론적 분류와 운영화 모델의 매핑 관계는
본문 \secref{app:operationalization}(제5장 실무적 기여)의 조작화 예시에서 참고용으로 다룬다.

\begin{figure}[htbp]
\centering
\begin{tikzpicture}[
  every node/.style={text=black},
  pillar/.style={draw, thick, rounded corners=4pt, fill=blue!10, minimum width=40mm,
    minimum height=22mm, font=\Large\bfseries, align=center, text=black},
  obs/.style={draw, very thick, rounded corners=4pt, fill=green!15, minimum width=64mm,
    minimum height=18mm, font=\Large\bfseries, align=center, text=black},
  arr/.style={->, >=Stealth, very thick, black}
]
  \node[pillar] (m) at (0, 0)   {Metrics\\\normalsize\mdseries(메트릭)};
  \node[pillar] (l) at (5.5, 0) {Logs\\\normalsize\mdseries(로그)};
  \node[pillar] (t) at (11, 0)  {Traces\\\normalsize\mdseries(트레이스)};
  \node[obs]    (o) at (5.5, -3.4)
    {Observability\\\normalsize\mdseries(관측성)};
  \draw[arr] (m.south) -- (o.north);
  \draw[arr] (l.south) -- (o.north);
  \draw[arr] (t.south) -- (o.north);
\end{tikzpicture}
\caption{관측성의 세 가지 핵심 구성 요소(three pillars)}
\label{fig:obs-three-pillars}
\end{figure}

관측성 수준이 높은 조직은 장애 발생 시 원인 파악에 소요되는 시간(Mean Time to Detect,
MTTD)과 복구에 소요되는 시간(Mean Time to Recovery, MTTR)을 모두 단축할 수 있다~\citep{forsgren2018accelerate, kim2016devops}.
이는 관측성이 단순한 기술적 역량을 넘어 조직의 복구 목표 달성에 직접적으로
기여하는 운영 변수임을 시사한다.
본 연구에서는 이러한 이론적 근거를 바탕으로 관측성 점수 $O_{\text{raw}} \in [0,1]$을
$\DRTO$ 모델의 핵심 입력 변수로 도입한다.


%% -------------------------------------------------------------------
\subsection{모니터링과 관측성의 개념적 구분}
\label{subsec:monitoring-vs-observability}
%% -------------------------------------------------------------------

실무에서 모니터링(monitoring)과 관측성(observability)은 혼용되는 경우가 많으나,
BCM 맥락에서 양자는 분석 범위와 목적이 구별된다. 모니터링은 사전에 정의된
지표를 수집하여 임계값 초과 여부를 감시하는 활동인 반면, 관측성은 시스템 외부
출력만으로 내부 상태를 추정할 수 있는 시스템 고유의 속성이다. 즉 모니터링이
관측성을 구현하기 위한 수단(도구와 프로세스)이라면, 관측성은 그 수단이 확보하는
결과(시스템의 파악 가능 정도)에 해당한다.
\p{[}표~\ref{tab:monitoring-vs-observability}\p{]}은 두 개념의 핵심 차이를 비교한 것이다.

\begin{table}[htbp]
\centering
\caption{모니터링과 관측성의 비교}
\label{tab:monitoring-vs-observability}
\small
\begin{tabularx}{\textwidth}{l X X}
\toprule
\textbf{비교 기준} & \textbf{모니터링 (Monitoring)} & \textbf{관측성 (Observability)} \\
\midrule
정의 &
  사전에 정의된 지표(메트릭, 로그, 알림 규칙)를 수집·추적하는 활동 &
  시스템 외부 출력만으로 내부 상태를 추정할 수 있는 시스템 고유 속성 \\
\addlinespace
관점 &
  ``무엇을 측정할 것인가'' (도구·프로세스 중심) &
  ``시스템이 얼마나 파악 가능한가'' (시스템 능력 중심) \\
\addlinespace
범위 &
  사전 정의된 이상 징후 탐지 (known-unknowns) &
  예상하지 못한 장애의 근본 원인 추적 포함 (unknown-unknowns) \\
\addlinespace
BCM 역할 &
  장애 감지 $\to$ 알림 발생 $\to$ 대응 절차 개시 &
  복구 과정에서 시스템 상태 가시성 확보 $\to$ 복구 시간 단축 \\
\addlinespace
$\DRTO$ 대응 &
  모니터링 도구의 존재 자체는 $O_{\mathrm{raw}}$의 필요조건 &
  $O_{\mathrm{raw}} \in [0,1]$로 정량화, 관측성 채널($E_{O,\mathrm{bnd}}$)에 직접 반영 \\
\bottomrule
\end{tabularx}
\end{table}

$\DRTO$ 모델에서는 모니터링 인프라의 수준이 $O_{\mathrm{raw}}$ 값으로 변환되어
관측성 채널에 반영되므로, 모니터링 투자가 곧 $\DRTO$ 감소로 이어지는 정량적
경로를 제공한다. 이러한 관측성의 차이는 실무 환경에서도 뚜렷하게 나타난다.
관측성이 높은($O_{\mathrm{raw}} \approx 1$) 환경의 전형적 예는
실시간 서버 모니터링 시스템(Prometheus, Grafana)이 구축된 IT 인프라로,
장애 발생 시 원인을 수 분 내에 특정하여 복구 시간을 단축하며,
금융권의 실시간 트랜잭션 모니터링이나 SRE 조직의 가시성 대시보드가 여기에 해당한다.
반대로 관측성이 낮은($O_{\mathrm{raw}} \approx 0$) 환경에서는
센서가 미비한 노후 제조 설비나 수동 점검에 의존하는 중소기업 IT 환경처럼
장비 고장의 근본 원인 파악에 수 시간이 소요되어 복구가 지연된다.

\chg{이처럼 관측성은 복구 환경의 가시성을 좌우하여 복구 시간에 직접 영향을 미치지만, 복구 시간에 내재한 변동성 자체를 설명하지는 못한다. 다음 절에서는 이 변동성을 정량화하는 불확실성 모델링을 검토한다.}
  % 2.3 관측성(Observability) 이론


%% ===================================================================
\section{불확실성(Uncertainty) 모델링}
\label{sec:lit-uncertainty}
%% ===================================================================

\chg{관측성이 복구 환경의 가시성을 높여 주더라도, 복구 소요 시간 자체에 내재한 확률적 변동까지 제거하지는 못한다. 이 절에서는 이러한 변동을 설명하는 불확실성 모델링을 일상적 변동을 다루는 가우시안 분포와 극단 사건을 다루는 파레토 분포를 중심으로 검토한다.}

%% -------------------------------------------------------------------
\subsection{가우시안 불확실성 분포}
\label{subsec:gaussian-uncertainty}
%% -------------------------------------------------------------------

불확실성 정량화의 가장 기본적인 도구는 가우시안(정규) 분포이다.
\citet{jaynes2003probability}는 확률 이론의 논리적 기초에 관한 체계적 논의를 통해
중심극한정리(Central Limit Theorem, CLT)의 정보이론적 정당성을 제시하였다.
다수의 독립적인 확률 변수의 합이 가우시안 분포에 수렴한다는 CLT는,
실제 복구 과정에서 발생하는 다양한 소규모 불확실성 요인, \chg{곧 인력 가용성과
장비 상태, 네트워크 지연, 절차적 오류 등의} 복합적 효과가
가우시안 형태로 모델링될 수 있음을 이론적으로 뒷받침한다.

본 연구에서는 일상적 운영 불확실성을 가우시안 분포 \chg{$U^{(G)} \sim N(\mu, \sigma^2)$}로
모델링한다. 이 분포는 대칭적 특성을 가지므로, 복구 시간의 단축과 연장 가능성을
동등하게 반영할 수 있다. 구체적으로 $\DRTO$ 프레임워크에서는
$U_{\text{raw}} \sim N(3, 1)$의 원시 불확실성 값을 생성한 후
$z_U = \kappa \cdot k_U \cdot R_0 \cdot U_{\text{raw}} / (R_{\max} - R_0)$으로
시그모이드 입력값을 산출한다.
\chg{가우시안을 C2 조건에 채택한 근거는 앞서 언급한 중심극한정리에 있다. 복구 과정에서
작용하는 다수의 독립적 소규모 불확실성 요인이 누적될 때 그 합이 가우시안에 수렴하므로,
어느 하나의 지배적 요인 없이 여러 작은 변동이 합산되는 일상적 운영 상황을 가우시안으로
표현하는 것이 이론적으로 타당하다.}
그러나 가우시안 분포는 꼬리 확률(tail probability)이 지수적으로 감소하므로,
대규모 중단이나 극단적 지연과 같은 드문 사건(rare event)의 발생 가능성을
구조적으로 과소평가하는 한계를 지닌다.


%% -------------------------------------------------------------------
\subsection{파레토 분포}
\label{subsec:heavy-tail}
%% -------------------------------------------------------------------

현실 세계의 재해 및 중단 사태는 가우시안 분포가 예측하는 것보다 훨씬 빈번하게
극단적 양상을 보인다. \citet{taleb2007blackswan}은 ``블랙 스완''의 개념을 통해
극도로 낮은 확률이지만 발생 시 막대한 영향을 미치는 사건들이 금융, 기술, 자연재해
영역에서 체계적으로 과소평가되고 있음을 경고하였다.

멱법칙(power-law) 분포는 이러한 극단적 사건을 수학적으로 표현하는 핵심 도구이다.
\citet{newman2005powerlaws}은 자연과학, 사회과학, 기술 시스템에서 발견되는
멱법칙 분포의 보편성을 체계적으로 검토하고, 그 확률밀도함수가
\begin{equation}
  f(x) = \frac{\alpha \, x_m^{\alpha}}{x^{\alpha+1}}, \quad x \geq x_m
  \label{eq:pareto-pdf}
\end{equation}
의 형태를 가지는 파레토(Pareto) 분포로 대표됨을 보였다.
여기서 $\alpha$는 형상 모수(shape parameter) 또는 꼬리 지수(tail index)로서,
분포의 꼬리 두께와 모멘트 존재 조건을 결정하는 핵심 파라미터이다.
$\alpha$가 작을수록 꼬리가 두꺼워져 극단값의 출현 빈도가 높아진다.

%% --- 형상 모수 α의 역할 ---
\subsubsection{형상 모수 $\alpha$의 역할}
\label{subsubsec:pareto-alpha}

파레토 분포의 통계적 성질은 형상 모수 $\alpha$의 값에 따라 질적으로 달라진다.
$n$차 모멘트 $E[X^n]$이 유한하기 위해서는 $\alpha > n$이어야 하며,
\chg{이로부터 다음과 같은 임계 구간이 도출된다. $\alpha$가 1 이하이면 평균조차 발산하여
모든 모멘트가 무한하고, 1과 2 사이이면 평균은 정의되나 분산이 발산한다. 2와 3 사이이면
평균과 분산은 유한하나 왜도가 발산하여 분포 형태가 극도로 비대칭이 되고, 3과 4 사이이면
왜도까지 유한하나 첨도가 발산하여 정규분포 대비 극단값의 빈도가 현저히 높아진다.
$\alpha$가 4를 넘으면 4차 이하의 모든 모멘트가 유한해져 꼬리가 상대적으로 얇아지면서
가우시안에 근사한다.}
\p{[}표~\ref{tab:pareto-alpha-properties}\p{]}는 대표적 $\alpha$ 값에 따른
파레토 분포($x_m = 1$)의 통계적 성질을 정리한 것이다.

\begin{table}[htbp]
  \centering
  \caption{형상 모수 $\alpha$에 따른 파레토 분포의 통계적 성질 ($x_m = 1$)}
  \label{tab:pareto-alpha-properties}
  \small
  \begin{tabular}{ccccc}
    \toprule
    $\alpha$ & 평균 $E[X]$ & 분산 $\text{Var}[X]$ & 왜도 & 첨도 \\
    \midrule
    1 & $\infty$ & $\infty$ & --- & --- \\
    2 & 2.000 & $\infty$ & $\infty$ & --- \\
    \textbf{3} & \textbf{1.500} & \textbf{0.750} & $\boldsymbol{\infty}$ & $\boldsymbol{\infty}$ \\
    5 & 1.250 & 0.104 & 4.649 & 70.800 \\
    10 & 1.111 & 0.015 & 2.811 & 14.828 \\
    \bottomrule
  \end{tabular}
\end{table}

$\alpha$가 증가할수록 평균은 $x_m$에 수렴하고 분산은 급격히 감소하여
분포가 $x_m$ 부근에 집중된다.
반대로 $\alpha$가 감소하면 꼬리가 두꺼워져
극단적으로 큰 값의 출현 확률이 높아진다.
이러한 $\alpha$의 역할은 생존함수(survival function)
$\bar{F}(x) = (x_m/x)^{\alpha}$를 통해 직관적으로 확인된다\chg{. 예컨대 $\alpha = 2$일 때
$\bar{F}(10) = 0.01$인 반면, $\alpha = 5$에서는 $\bar{F}(10) = 10^{-5}$로 극단값의 출현
확률이 1/1{,}000 이하로 감소한다.}

사이버 위험(cyber risk) 분야에서 멱법칙 분포의 적합성은 다수의 실증 연구에 의해
확인되어 왔다. 데이터 유출 규모, 사이버공격에 의한 재무 손실, 시스템 장애 지속 시간 등은
정규분포보다 멱법칙 분포에 의해 더 잘 설명되며,
이는 BCM 맥락에서 극단적 복구 지연 시나리오의 모델링에
파레토 분포가 적합함을 시사한다.
이에 본 연구에서는 멱법칙 분포의 대표적 형태인 파레토 분포를 채택하여
C3 조건의 불확실성을 모형화하며,
이후 본문에서 ``파레토 분포''로 통일 표기한다.

극단값 이론(Extreme Value Theory, EVT)은 이러한 파레토 분포의 극단적 사건을
체계적으로 분석하는 통계적 프레임워크를 제공한다~\citep{coles2001extremes}.
\citet{embrechts1997extremal}은 보험 및 금융 분야에서의 극단값 모델링을 위한
수학적 기초를 수립하였으며, \citet{mcneil2015risk}는 정량적 위험관리에서
EVT의 활용 방안을 구체적으로 논의하였다.
\chg{본 연구는 극단값 이론을 직접적인 추정 기법으로 사용하지는 않으나, 그 핵심 결과는
C3 조건에서 파레토 분포를 채택한 이론적 정당성을 제공한다. 극단값 이론에 따르면 광범위한
분포에서 충분히 높은 임계값을 초과하는 부분이 일반화 파레토 분포로 수렴하므로, 복구 시간의
극단적 꼬리를 파레토 분포로 모형화하는 것은 임의의 선택이 아니라 극단값의 점근적 거동에
부합하는 선택이다. 본 연구가 멱법칙과 극단값 이론에서 가져온 것은 ``극단적 사건의 꼬리는
지수적으로가 아니라 멱법칙으로 감소한다''는 명제이며, 이를 C3 채널의 분포 형태로 직접
반영한 것이다.}

%% --- α = 3 선택 근거 ---
\subsubsection{$\alpha = 3$ 선택 근거}
\label{subsubsec:alpha3-rationale}

본 연구에서는 $\DRTO$ 프레임워크의 조건 \textbf{C3}에서 파레토 분포
\chg{$U^{(P)} \sim 6 \cdot \text{Pareto}(\alpha = 3)$}에 기반한 불확실성 채널을 도입하여,
가우시안 분포(\textbf{C2})가 포착하지 못하는 극단적 복구 지연 시나리오를 모델링한다.
\chg{$\alpha = 3$의 선택은 네 가지 근거에 기반한다. 첫째, $\alpha > 2$이면 평균과 분산이
모두 유한하여 표본 통계량의 안정적 추정과 회귀분석이 가능하다. 둘째, $\alpha \leq 4$이면
첨도가 발산하여 가우시안 분포 대비 극단값이 현저히 빈번하게 출현하는 현실적 시나리오를
표현한다. 셋째, $\alpha = 3$은 유한 분산과 무한 왜도 및 첨도가 공존하는 경계에 위치하여
통계적 분석 가능성과 중꼬리(heavy-tail) 특성 사이의 최적 균형점을 제공한다. 넷째,
사이버 보안 사고의 규모 분포에서 보고되는 $\alpha$ 추정값이 2와 4 사이의 중앙 영역에
해당한다는 실증적 정합성을 갖는다.}


%% --- C2·C3 정량적 비교 및 첨도-꼬리 관계 ---
\subsubsection{가우시안과 파레토 분포의 정량적 비교}
\label{subsubsec:c2c3-comparison}

$\DRTO$ 프레임워크가 채택한 두 불확실성 채널, 즉 가우시안 분포
\chg{$U^{(G)} \sim N(3,\,1)$}(C2)과 파레토 분포 \chg{$U^{(P)} \sim 6 \cdot \text{Pareto}(\alpha=3)$}(C3)은
동일한 기대값($3.0$)을 갖도록 척도를 조정하였으나, 꼬리 영역의 거동에서 질적으로
구별된다. \p{[}표~\ref{tab:c2c3-comparison}\p{]}은 두 분포의 꼬리 성질과 주요 분위수를
비교한 것이다.

\begin{table}[htbp]
\centering
\caption{가우시안(C2)과 파레토(C3) 불확실성 분포의 비교}
\label{tab:c2c3-comparison}
\small
\begin{tabular}{l c c l}
\toprule
\textbf{구분} & \textbf{가우시안 (C2)} & \textbf{파레토 (C3)} & \textbf{의미} \\
\midrule
꼬리 분류 & 경꼬리 (light-tail) & 중꼬리 (heavy-tail) & 감소 속도 기반 성질 분류 \\
꼬리 감소 속도 & $\sim e^{-x^2/2}$ (지수적) & $\sim x^{-\alpha}$ (멱법칙) & 중꼬리가 훨씬 느리게 감소 \\
첨도 & $3$ (기준값) & $\infty$ ($\alpha = 3$) & 꼬리 두께의 수치 지표 \\
분산 & $1.0$ & $27.0$ (약 $27$배) & 변동 폭의 차이 \\
왜도 & $0$ (대칭) & 양수 (우편향) & 극단 지연의 비대칭성 \\
P95 & $4.645$ & $10.29$ (약 $2.2$배) & 상위 $5\%$ 복구 시간 \\
P99 & $5.326$ & $21.85$ (약 $4.1$배) & 상위 $1\%$ 복구 시간 \\
중앙값 & $\approx 3.0$ & $1.56$ & 대부분은 빠르되 소수가 급격히 지연 \\
\bottomrule
\end{tabular}
\end{table}

여기서 첨도와 꼬리 분류는 관련은 있으나 서로 다른 개념임에 유의해야 한다.
첨도는 분포의 꼬리 두께를 하나의 수치로 측정한 지표인 반면, 경꼬리와 중꼬리는
꼬리의 감소 속도에 따른 성질 분류이다. 중꼬리 분포는 첨도가 높은 경향이 있으나,
첨도가 높다고 반드시 중꼬리인 것은 아니다. 예컨대 유한 지지(bounded support)를
갖는 분포도 극단적 이봉 구조를 통해 높은 첨도를 가질 수 있으나 경꼬리에 속한다.
파레토($\alpha=3$)는 첨도가 무한대이므로 중꼬리의 전형적 사례에 해당한다.
파레토 분포의 중앙값($1.56$)이 평균($3.0$)보다 현저히 낮다는 사실은, 대부분의 사건이
빠르게 복구되지만 소수의 극단 사례에서 복구가 급격히 지연되어 평균을 끌어올리는
중꼬리의 구조적 특성을 직접 보여준다. 두 분포의 시각적 비교와 P99 수치는
\figref{fig:pareto_concept}(본문 \secref{sec:result_c3})를 함께 참조하라.

이러한 분포 특성의 차이는 실무 맥락에서도 뚜렷이 구별된다. 가우시안 불확실성은
일상적 운영 변동을 대표하는데, 일반적 라우터 장애의 복구 시간이나 정기적 소프트웨어
업데이트의 롤백, 표준화된 서버 교체 절차처럼 복구 시간이 평균 근처에 집중되고
극단적 편차의 확률이 지수적으로 감소하여 대부분 예측 가능한 범위에서 마무리되는
경우가 여기에 해당한다. 경꼬리 특성상 이러한 사건은 P99에서도 평균 대비
$2.3\sigma$ 이내로 한정된다. 반면 파레토 불확실성은 극단적 사건을 대표하는데,
전사 시스템이 암호화되는 랜섬웨어 공격에서 대부분은 백업 복원으로 수 시간 내에
해결되나 소수의 극단 사례에서는 데이터 복호화와 포렌식 분석으로 복구가 수일까지
연장되며(P99에서 약 $13.2$시간 이상), 지진 후 화재와 정전이 동시에 발생하는 복합
재난이나 핵심 부품 공급사의 공장 화재로 인한 공급망 연쇄 장애 또한 개별 사건
복구 시간의 합을 초과하는 비선형적 지연을 유발하여 꼬리 영역의 복구 부담을 급증시킨다.


%% -------------------------------------------------------------------
\subsection{불확실성 분류 체계의 통합적 이해}
\label{subsec:uncertainty-classification}
%% -------------------------------------------------------------------

본 절에서 다룬 가우시안 분포(C2)와 파레토 분포(C3)는 불확실성을 \textbf{분포 형태}
(즉 \textbf{강도 분포}의 모양) 차원에서 분류한 것이다. 그러나 BCM 운영 환경에서 불확실성 전반을 체계적으로
이해하기 위해서는 \citet{walker2003uncertainty}이 제시한 \textbf{불확실성 단계 분류}
프레임워크와 함께 살펴볼 필요가 있다(\p{[}그림~\ref{fig:uncertainty-levels}\p{]}).
이 프레임워크는 우리가 알고 있는 정도(level of knowledge)에 따라 불확실성을 네 단계로
구분한다. 이는 리스크 평가의 표준 분해(리스크 = 빈도 × 강도)와 대응되며, 분포는 강도
(Severity)에, 시나리오·확률은 빈도(Frequency)에 해당한다. \chg{네 단계는 i) 분포와
확률을 모두 추정할 수 있는 통계적 불확실성, ii) 대안 시나리오는 열거할 수 있으나 확률이
미정인 시나리오 불확실성, iii) 분포 형태 자체를 특정할 수 없는 인지된 무지, iv) 모델과
시나리오와 확률이 모두 불확실한 전적 무지로 구분된다.}

본 연구의 $\DRTO$ 모델은 이 분류 체계 중 양 극단을 정량적으로 다룬다.
통계적 불확실성은 가우시안 분포(C2, \chg{$U^{(G)} \sim N(3,1)$})에 대응하며, 인지된 무지는
파레토 분포(C3, \chg{$U^{(P)} \sim 6 \cdot \text{Pareto}(\alpha\!=\!3)$})로 보수적으로 표현된다.
시나리오 불확실성은 관측성 채널($O_{\text{raw}}$, $k_O$)의 정보 확보를 통해 부분적으로
흡수되며, 전적 무지(블랙 스완 사건)는 본 모델의 적용 한계로 4.3절에서 별도 논의한다.
\chg{한편 운영 단계에서 $U_{\text{raw}}$를 추정하기 위한 데이터 출처(과거 장애 이력 DB,
외부 위협 지수)와 합성 방식은 본문 \subsecref{app:operationalization} 조작화 예시에서 참고용으로 다룬다.}

\begin{figure}[htbp]
\centering
\begin{tikzpicture}[
  every node/.style={text=black},
  level/.style={draw, thick, rounded corners=3pt, text width=125mm,
    minimum height=24mm, font=\large\bfseries, align=left, inner sep=6pt, text=black},
  arr/.style={-{Stealth[length=5mm]}, ultra thick, black}
]
  \node[level, fill=green!15]                   (l1) {Statistical uncertainty (통계적 불확실성)\\
    \normalsize\mdseries 분포(강도)·확률(빈도) 추정 가능. 반복 관측으로 정량화. 예: 장비 고장률.
    \;\textit{$\to$ DRTO C2 가우시안}};
  \node[level, fill=blue!12, below=2mm of l1]   (l2) {Scenario uncertainty (시나리오 불확실성)\\
    \normalsize\mdseries 대안 시나리오 열거 가능, 확률(빈도) 미정. 예: 미경험 재난 유형.
    \;\textit{$\to$ 관측성 채널}};
  \node[level, fill=orange!18, below=2mm of l2] (l3) {Recognized ignorance (인지된 무지)\\
    \normalsize\mdseries 분포(강도) 형태 자체 특정 불가. 예: 신종 사이버공격, 블랙 스완.
    \;\textit{$\to$ DRTO C3 파레토}};
  \node[level, fill=red!12, below=2mm of l3]    (l4) {Total ignorance (전적 무지)\\
    \normalsize\mdseries 모델·시나리오·확률(빈도·강도) 모두 불확실. 예: 기후 복합재난.
    \;\textit{$\to$ 4.3절 한계}};

  \coordinate (atop) at ($(l1.north west) + (-12mm, 0)$);
  \coordinate (abot) at ($(l4.south -| atop)$);
  \draw[arr] (atop) -- (abot);
  \node[font=\normalsize\bfseries, rotate=90, text=black, anchor=south]
    at ($(atop)!0.5!(abot) + (-3mm, 0)$) {불확실성 정도 증가};
\end{tikzpicture}
\caption{불확실성의 단계별 분류 체계와 $\DRTO$ 모델 대응}
\label{fig:uncertainty-levels}
\end{figure}


%% -------------------------------------------------------------------
\subsection{불확실성 하위개념과 업무연속성 리스크와의 관계}
\label{subsec:uncertainty-risk}
%% -------------------------------------------------------------------

\chg{본 연구에서 불확실성은 복구시간의 확률적 변동을 지칭하며, BCM에서
다루는 리스크(재난)와는 분석 수준이 다르다. 리스크는 무엇이 발생하는가라는
사건(event)에 초점을 두는 반면, $\DRTO$에서의 불확실성은 복구시간이 얼마나 변동하는가라는
결과(outcome)에 초점을 둔다. 즉 리스크 사건이 발생하면 복구시간에 불확실성이
유입되고, $\DRTO$가 이를 확률분포로 정량화한다.

경제학자 \citet{knight1921risk}는 리스크와 불확실성을 다음과 같이 구분하였다. 리스크는 결과의
확률분포가 알려져 정량적 측정이 가능한 상황을, 불확실성은 확률분포 자체를 특정할 수
없는 상황을 의미한다. 이 구분은 본 연구의 분류 체계와 대응한다. C2의 가우시안
불확실성은 분포 형태와 매개변수가 사전에 특정되므로 나이트적 의미의 리스크에 가깝고,
C3의 파레토 불확실성은 분포 형태를 확정할 수 없는 나이트적 불확실성을 보수적으로
근사한 것이다. 불확실성을 학술적으로 더 세분한 분류와 본 모델의 대응은
\p{[}표~\ref{tab:app-uncertainty-taxonomy}\p{]}과 같다.}

\begin{table}[htbp]
\centering
\caption{불확실성 하위개념 분류와 $\DRTO$ 모델 대응}
\label{tab:app-uncertainty-taxonomy}
\small
\begin{tabular}{p{1.8cm} p{2.0cm} p{3.8cm} p{2.8cm} p{2.8cm}}
\toprule
\textbf{분류} & \textbf{영어 표기} & \textbf{정의} & \textbf{BCM 예시} & \textbf{DRTO 대응} \\
\midrule
인식론적 불확실성 & Epistemic & 지식 부족에 기인, 정보 확보로 축소 가능 & 미경험 재난 유형, 불완전한 BIA & 관측성 채널 ($O_{\text{raw}}$, $k_O$) \\
\midrule
우연적 불확실성 & Aleatory & 본질적 무작위성, 데이터를 추가해도 축소 불가 & 장비 고장률, 네트워크 지연 변동 & C2 가우시안 ($U^{(G)}$) \\
\midrule
나이트적 불확실성 & Knightian & 확률분포 자체를 특정할 수 없는 불확실성 & 블랙 스완 사건, 신종 사이버공격 & C3 파레토 ($U^{(P)}$) \\
\midrule
심층 불확실성 & Deep Uncertainty & 모델·시나리오·확률 모두 불확실 & 기후변화 복합재난, 팬데믹 & 한계 사항 (\secref{sec:limitations}) \\
\bottomrule
\end{tabular}
\end{table}

\chg{표에서 우연적 불확실성은 반복 관측으로 분포를 추정할 수 있으나 변동 자체를 제거할
수는 없는 유형이고, 인식론적 불확실성은 관측성($O_{\text{raw}}$)을 높임으로써 축소할 수
있어 $\DRTO$ 모델의 관측성 채널이 직접 반영하는 메커니즘에 해당한다. 나이트적 불확실성은
분포 형태를 사전에 특정할 수 없으므로 본 연구에서는 파레토 분포로 보수적으로 모델링하며,
심층 불확실성은 현 모델의 범위를 벗어나 \secref{sec:limitations}에서 한계로 명시하였다.}

\chg{한편 BCM에서 다루는 세 가지 재난 유형과 본 모델의 불확실성은
\p{[}표~\ref{tab:app-disaster-mapping}\p{]}과 같이 대응한다. 동일한 재난 유형이라도
발생 규모에 따라 불확실성 특성이 달라져, 일상적 서버 장애는 가우시안(C2)으로,
대규모 사이버공격은 파레토(C3)로 모형화된다.}

\begin{table}[htbp]
\centering
\caption{BCMS 재난 유형과 $\DRTO$ 불확실성 매핑}
\label{tab:app-disaster-mapping}
\small
\begin{tabular}{p{2cm} p{3.2cm} p{2.8cm} p{2.5cm} p{2.5cm}}
\toprule
\textbf{재난 유형} & \textbf{세부 예시} & \textbf{빈도-영향 특성} & \textbf{불확실성 유형} & \textbf{DRTO 조건} \\
\midrule
자연재난 & 지진, 태풍, 홍수, 산사태 & 빈도 낮음, 영향 극대 & 나이트적 & C3 (파레토) \\
\midrule
사회재난 & 대규모 정전, 통신장애, 공급망 단절, 테러 & 빈도 중간, 영향 중$\sim$대 & 우연적 + 나이트적 & C2 $\sim$C3 \\
\midrule
기술적 재난 & 서버 장애, 사이버공격, 소프트웨어 결함, 데이터 유출 & 빈도 높음, 영향 소$\sim$극대 & 우연적(일상) / 나이트적(극단) & C2(일상) / C3(극단) \\
\bottomrule
\end{tabular}
\end{table}

\chg{사회재난의 경우에도 공급망 단절의 파급 범위에 따라 가우시안(C2)에서 파레토(C3)로
전이될 수 있다.}

\chg{끝으로 BCM에서의 리스크와 $\DRTO$에서의 불확실성의 유사점과 차이점을
\p{[}표~\ref{tab:app-risk-vs-uncertainty}\p{]}에 비교하였다. 리스크가 사건의 발생
가능성과 영향을 다루는 데 비해 불확실성은 그 결과인 복구시간의 확률적 변동을 다루며,
양자는 모두 복구시간에 영향을 미치므로 $\DRTO$가 이를 통합 정량화한다.}

\begin{table}[htbp]
\centering
\caption{BCMS 리스크(재난)와 $\DRTO$ 불확실성 비교}
\label{tab:app-risk-vs-uncertainty}
\small
\begin{tabular}{p{2.5cm} p{5.5cm} p{5.5cm}}
\toprule
\textbf{비교 기준} & \textbf{BCMS 리스크 (재난)} & \textbf{DRTO 불확실성} \\
\midrule
정의 & 조직 목표에 부정적 영향을 미치는 사건의 발생 가능성과 결과~\nrev{(ISO 31000, 2018)} & 복구시간($\DRTO$)의 확률적 변동 범위 \\
\midrule
분석 단위 & 사건(event): ``무엇이 발생하는가'' & 결과(outcome): ``복구시간이 얼마나 변동하는가'' \\
\midrule
방향성 & 주로 부정적 (위협 중심) & 양방향 (단축도 연장도 가능) \\
\midrule
정량화 방법 & 발생가능성 $\times$ 영향도 매트릭스 & 확률분포 (가우시안 / 파레토) \\
\midrule
시간적 성격 & BIA 시점 고정 평가 (정적) & 실시간 동적 변동 \\
\midrule
축소 방법 & 예방·완화 통제 수단 & 관측성 향상 ($O_{\text{raw}}$ $\uparrow$) \\
\midrule
\multicolumn{3}{l}{공통점: 둘 다 복구시간에 영향 $\to$ $\DRTO$가 양자를 통합 정량화} \\
\bottomrule
\end{tabular}
\end{table}

\chg{이상에서 살펴본 가우시안과 파레토 불확실성은 모두 복구 시간에 확률적 변동을 유발하므로, 이를 $\DRTO$ 모델에 반영하려면 그 영향을 일정한 경계 안에서 변환하는 장치가 필요하다. 다음 절에서는 이러한 변환을 담당하는 시그모이드 함수와 바운딩 기법을 검토한다.}
  % 2.4 불확실성(Uncertainty) 모델링


%% ===================================================================
\section{시그모이드 함수와 바운딩}
\label{sec:lit-sigmoid}
%% ===================================================================

%% -------------------------------------------------------------------
\subsection{시그모이드 함수의 수학적 성질}
\label{subsec:sigmoid-properties}
%% -------------------------------------------------------------------

$\DRTO$ 모델에서는 관측성과 불확실성의 복합 효과를 물리적으로 타당한 범위 내로
제한하는 바운딩(bounding) 기법이 필수적이다. 복구 시간은 물리적으로
음수가 될 수 없으며, 조직이 허용하는 최대 복구 시간($R_{\max}$)을 초과할 수도 없기 때문이다.
기준 RTO($R_0$)와 최대 복구 시간 MTPD($R_{\max}$)의 상세 정의는 2.7절에서 제시한다.

\chg{본 연구의 바운딩 함수는 \citet{verhulst1838notice}가 인구 성장 모형으로 처음
도입한 표준 로지스틱 함수 $\sigma(z) = 1/(1+e^{-z})$에 이론적 뿌리를 둔다. 로지스틱
함수는 매끄러운 S자 형태로 출력을 유한 구간 안에 제한하는 성질 덕분에, 생물학적 성장
모형에서 출발하여 통계학과 여러 응용 분야의 표준적 변환으로 정착해 왔다~\citep{cramer2002origins}.}
본 연구에서 채택한 수정 시그모이드(modified sigmoid) 함수는 다음과 같이 정의된다.
\begin{equation}
  S(z) = \frac{2}{1 + e^{-z}} - 1, \quad S: \mathbb{R} \to (-1,\; 1)
  \label{eq:sigmoid-def}
\end{equation}

이 함수는 표준 로지스틱 함수를 $(-1, +1)$ 구간으로 사상한 것이다.
\chg{표준 로지스틱 함수의 치역이 $(0,1)$인 것과 달리 이를 $(-1,+1)$로 사상하여 채택한
이유는, $\DRTO$ 모델에서 관측성은 복구 시간을 기준선 \emph{아래}로, 불확실성은 기준선
\emph{위}로 작용하므로 $S(0)=0$을 중심으로 음과 양의 양방향으로 대칭인 출력이 필요하기
때문이다. 즉 로지스틱 함수의 ``매끄러운 유한 바운딩''이라는 성질은 그대로 가져오되,
두 채널의 상하 대칭 작용을 표현하기 위해 치역만 원점 대칭으로 \chg{평행이동하고 확대한} 것이다.}
%% [G19/일원화] 구 성질 열거(S(0)=0·극한·도함수·홀함수·C∞)는 2.9절 subsec:basic-properties(2계 도함수·변곡점 포함 엄밀 정리)와 중복 → 한 문장으로 축약·전방참조.
\chg{이 함수는 $S(0) = 0$의 중립성, $(-1, +1)$로의 매끄러운 유한 바운딩, 전 구간 단조증가,
$S(-z) = -S(z)$의 원점 대칭(홀함수), $C^{\infty}$ 미분 가능성 등의 핵심 성질을 가지며, 이들
성질의 엄밀한 정리와 2계 도함수와 변곡점 분석은 모델의 수학적 성질을 다루는 2.9절에서 종합한다.}


%% -------------------------------------------------------------------
\subsection{\chg{선형과 클리핑, 시그모이드 바운딩 기법의 비교}}
\label{subsec:bounding-comparison}
%% -------------------------------------------------------------------

\chg{$\DRTO$ 모델의 바운딩 함수는 다음 다섯 가지 설계 요구조건을 충족해야 한다. 첫째,
출력이 언제나 물리적으로 타당한 범위 $0 < \DRTO < R_{\max}$ 안에 머물러야 한다(경계 보장).
둘째, 입력이 작은 평균 영역에서는 근사적으로 선형으로 반응하여 해석이 직관적이어야 한다
(평균 영역 선형성). 셋째, 입력이 극단적으로 커지는 영역에서는 출력이 포화하여 경계를
넘지 않아야 한다(극단 영역 포화). 넷째, 관측성과 불확실성 두 채널의 기여가 각각의 물리적
범위 안에서 분리되어 작용해야 한다(채널 분리). 다섯째, 입력 변화에 대해 출력이 단조적으로
반응하여 입력의 순서가 보존되어야 한다(단조성). 이 다섯 요구조건을 기준으로 세 가지 후보
바운딩 기법을 비교한다.}

\textbf{선형 모델}은 $\DRTO = R_0 + \alpha \cdot x$ 형태로서, 구현이 가장 단순하고
해석이 용이하다. 그러나 입력값이 극단적인 경우 $\DRTO$가 음수가 되거나
$R_{\max}$를 초과하는 물리적 비현실성이 발생할 수 있어 \chg{첫 번째 요구조건인 경계
보장을 충족하지 못한다}.

\textbf{클리핑(hard clipping)}은 $f(z) = \min(\max(z, z_{\min}), z_{\max})$로서
입력값이 경계 $[z_{\min}, z_{\max}]$ 내이면 그대로 통과시키고 경계를 초과하면 강제로
절단하는 경성 바운딩 함수이다. 클리핑은 경계 조건을 강제로 보장하나, 경계 지점에서
도함수가 불연속이 되고 \chg{경계 외의 다수 입력값이 단일 경계값으로 매핑되면서} 경계
근방의 미세한 입력 변화가 출력에 전혀 반영되지 않는 정보 손실이 발생한다.
\chg{즉 클리핑은 경계 보장 요구조건은 충족하지만, 경계 지점에서 매끄럽지 않아 평균 영역의
선형성과 극단 영역의 점진적 포화를 함께 만족하지 못하며, 경계 외에서 입력 변화가 출력에
반영되지 않아 단조성도 부분적으로 훼손된다.}

\textbf{시그모이드 바운딩}은 매끄러운 S자 곡선을 통해 출력을 유한 범위 내로
자연스럽게 제한하며, 전 구간에서 미분 가능하여 입력 변수의 변화에 대한
출력의 민감도를 연속적으로 분석할 수 있다.
\chg{시그모이드 바운딩은 매끄러운 유한 제한으로 경계를 보장하고, 원점 부근에서 근사적으로
선형이며, 입력이 커질수록 출력이 점진적으로 포화하고, 두 채널을 각각의 물리적 범위에서
독립적으로 바운딩할 수 있으며, 전 구간에서 단조증가한다. 따라서 앞서 제시한 다섯 가지
설계 요구조건을 모두 충족하는 것은 세 후보 중 시그모이드 바운딩뿐이며, 본 연구가 수정
시그모이드를 채택한 근거가 여기에 있다.}

\p{[}표~\ref{tab:bounding-comparison}\p{]}에 세 기법의 특성을 종합적으로 비교하였다.

\begin{table}[htbp]
\centering
\caption{바운딩 기법의 비교}
\label{tab:bounding-comparison}
\small
\begin{tabularx}{\textwidth}{@{}>{\raggedright\arraybackslash}p{2.4cm}
  >{\raggedright\arraybackslash}X
  >{\raggedright\arraybackslash}X
  >{\raggedright\arraybackslash}X@{}}
\toprule
\textbf{특성} & \textbf{선형 모델} & \textbf{클리핑(경성)} & \textbf{시그모이드 바운딩} \\
\midrule
수식 &
$R_0 + \alpha x$ &
$\min(\max(\cdot, 0), R_{\max})$ &
$R_0 + E_{O,\text{bnd}} + E_{U,\text{bnd}}$ \\
\addlinespace
경계 보장 &
미보장 (음수/초과 가능) &
보장 (강제 절단) &
보장 (점근적 접근) \\
\addlinespace
미분 가능성 &
전 구간 ($C^{\infty}$) &
경계점 불연속 &
전 구간 ($C^{\infty}$) \\
\addlinespace
경계 정보 보존 &
해당 없음 &
경계에서 기울기 $= 0$ (정보 손실) &
경계에서도 기울기 $> 0$ (정보 보존) \\
\addlinespace
포화 효과 &
없음 (무한 선형 증가) &
경성 포화 (즉시 절단) &
연성 포화 (점진적 둔화) \\
\addlinespace
물리적 해석 &
비현실적 극단값 가능 &
현실적이나 불연속 &
현실적이고 연속적 \\
\bottomrule
\end{tabularx}
\end{table}

\figref{fig:bounding-comparison}\refpo{fig:bounding-comparison}{은}{는} 세 바운딩 기법의 출력 특성을 비교한 것이다.
\chg{선형 모델은 경계를 초과할 수 있고 클리핑은 경계점에서 불연속인 반면, 시그모이드는 $(-1, +1)$ 범위 안에서 경계에 매끄럽게 점근하면서도 전 구간에서 기울기가 양으로 유지되는 점근적 바운딩을 제공한다.}
시그모이드 바운딩은 경계 영역에서의 정보 보존성, 미분 가능성,
운영 해석 가능성 측면에서 $\DRTO$ 프레임워크에 가장 적합하다.

\begin{figure}[htbp]
  \centering
  \begin{tikzpicture}
    \begin{axis}[
      width=0.88\textwidth,
      height=0.52\textwidth,
      xlabel={입력 $z$},
      ylabel={출력},
      xmin=-4.5, xmax=4.5,
      ymin=-1.6, ymax=1.9,
      grid=major,
      grid style={black},
      legend pos=south east,
      legend style={font=\small, rounded corners=2pt, fill=white, fill opacity=0.9},
      axis lines=middle,
      every axis x label/.style={at={(axis description cs:1.02,0.5)}, anchor=west, font=\small},
      every axis y label/.style={at={(axis description cs:0.5,1.02)}, anchor=south, font=\small},
      clip=false
    ]
      % 선형 모델 (제한 없음)
      \addplot[thick, blue!70, domain=-4.5:4.5, samples=100]
        {0.4*x};
      \addlegendentry{선형 ($0.4z$)}

      % 클리핑
      \addplot[thick, red!70, domain=-4.5:4.5, samples=200]
        {min(max(0.4*x, -1), 1)};
      \addlegendentry{클리핑 (hard)}

      % 시그모이드 S(z) = 2/(1+exp(-z)) - 1
      \addplot[thick, black, domain=-4.5:4.5, samples=200]
        {2/(1+exp(-x)) - 1};
      \addlegendentry{시그모이드 $S(z)$}

      % 경계 초과 영역 채우기 (위쪽: x>2.5에서 0.4x > 1)
      \fill[blue!12, opacity=0.5]
        (axis cs:2.5,1) -- (axis cs:4.5,1.8) -- (axis cs:4.5,1) -- cycle;

      % 경계 초과 영역 채우기 (아래쪽: x<-2.5에서 0.4x < -1)
      \fill[blue!12, opacity=0.5]
        (axis cs:-2.5,-1) -- (axis cs:-4.5,-1.8) -- (axis cs:-4.5,-1) -- cycle;

      % 경계선 표시
      \addplot[thin, dashed, black] coordinates {(-4.5, 1) (4.5, 1)};
      \addplot[thin, dashed, black] coordinates {(-4.5, -1) (4.5, -1)};

      % 경계 라벨
      \node[font=\scriptsize, text=black, anchor=west] at (axis cs:4.6, 1) {$+1$};
      \node[font=\scriptsize, text=black, anchor=west] at (axis cs:4.6, -1) {$-1$};

      % 선형 경계 초과 영역 라벨 (위쪽)
      \node[font=\scriptsize, text=black, fill=white, inner sep=1.5pt,
        anchor=south west] at (axis cs:2.8, 1.45)
        {선형: 경계 초과};
      \draw[-{Stealth[length=1.5mm]}, blue!70, thin]
        (axis cs:3.2, 1.42) -- (axis cs:3.5, 1.15);

      % 선형 경계 초과 영역 라벨 (아래쪽)
      \node[font=\scriptsize, text=black, fill=white, inner sep=1.5pt,
        anchor=north east] at (axis cs:-2.8, -1.45)
        {선형: 경계 초과};
      \draw[-{Stealth[length=1.5mm]}, blue!70, thin]
        (axis cs:-3.2, -1.42) -- (axis cs:-3.5, -1.15);

    \end{axis}
  \end{tikzpicture}
  \caption{바운딩 기법 비교}
  \label{fig:bounding-comparison}
\end{figure}


%% -------------------------------------------------------------------
\subsection{개별 시그모이드 바운딩의 적용}
\label{subsec:individual-sigmoid}
%% -------------------------------------------------------------------

본 연구에서 채택한 개별 시그모이드 바운딩(individual sigmoid bounding) 모델은
관측성과 불확실성의 효과를 각각 고유한 물리적 범위(span)에서
독립적으로 바운딩하는 구조이다. 대안인 결합(joint) 시그모이드 $S(z_O + z_U)$가
두 효과를 단일 함수로 압축하는 반면, 개별 시그모이드는 각 효과의 물리적 범위를
독립적으로 보장한다.
%% [G19/일원화] 구 핵심수식(eq:drto-individual-sigmoid, eq:lit-E_O_bnd, eq:lit-E_U_bnd, eq:lit-z_O, eq:lit-z_U)은
%% 2.7절 모델 구조(eq:drto, eq:E_O_bnd, eq:E_U_bnd, eq:z_O, eq:z_U; canonical)와 완전 중복 → 수식 삭제·전방참조(라벨 참조 0건 확인).
\chg{구체적으로 최종 $\DRTO$는 기준선 $R_0$에 관측성 바운딩 조정량
$E_{O,\text{bnd}} = -R_0 \cdot S(z_O)$(물리적 범위 $R_0$)와 불확실성 바운딩 조정량
$E_{U,\text{bnd}} = (R_{\max} - R_0) \cdot S(z_U)$(물리적 범위 $R_{\max} - R_0$)를 가산하여
산출한다. 음수 부호 $-R_0$는 관측성이 높을수록 복구 시간 단축 효과로 작용함을, $z_U$의
분모 $R_{\max} - R_0$는 불확실성 효과를 가용 상향 여유 범위로 정규화함을 뜻한다. 시그모이드
입력 인자 $z_O$, $z_U$의 전체 정식과 변수 및 매개변수의 형식적 정의, 경계 보장의 엄밀한
증명, 곡률 $\kappa = 1.2$의 선정 근거는 모델 구조와 최적화를 다루는 2.7절과 2.10절에서
종합하므로, 본 절은 그 형태만 바운딩 기법의 귀결로 제시한다.}

\chg{이 구조는 세 가지 핵심적 장점을 갖는다. 첫째, 경계 보장의 보편성으로서, 관측성 효과의
바운딩 범위($R_0$)와 불확실성 효과의 바운딩 범위($R_{\max} - R_0$)를 각각의 물리적 의미에
맞게 독립적으로 설정하므로 $\DRTO \in [0,\; R_{\max}]$가 $\kappa$ 값에 무관하게 수학적으로
보장된다. 둘째, 경사 최적화와 민감도 분석의 용이성으로서, 각 시그모이드 함수가 전체
정의역에서 $C^{\infty}$이므로 경사 기반 최적화와 민감도 분석이 용이하다. 셋째, C1 조건
분리 가능성으로서, C1 조건($E_{U,\text{bnd}} = 0$)에서 관측성만의 순수 효과를 불확실성
상한과 독립적으로 분석할 수 있어 모델의 단계적 확장이 자연스럽다.}


%% -------------------------------------------------------------------
\subsection{이중채널 감쇠 메커니즘}
\label{subsec:dual-channel-attenuation}
%% -------------------------------------------------------------------

개별 시그모이드 바운딩 구조는 관측성 채널($E_{O,\text{bnd}} = -R_0 \cdot S(z_O)$,
복구 시간 단축)과 불확실성 채널($E_{U,\text{bnd}} = (R_{\max}-R_0) \cdot S(z_U)$,
복구 시간 연장)이 동시에 작용하는 이중채널 구조를 형성한다. 두 채널이 각자의 물리적
범위에서 독립적으로 바운딩되므로, 불확실성이 극단적으로 커지는 꼬리(tail) 영역에서는
불확실성 채널이 지배적으로 작용하면서 관측성 채널의 단축 효과가 구조적으로 약화되는
현상이 나타난다. 이를 이중채널 감쇠(dual-channel attenuation)라 한다.

이 메커니즘은 실무적으로 중요한 함의를 갖는다. 고도의 모니터링 체계를 갖춘
조직(높은 관측성)이라 하더라도 대규모 사이버 공격과 같은 극단적 불확실성 상황에서는
관측성 투자만으로 복구 시간을 충분히 단축하기 어려우며, 불확실성 자체에 대한 관리,
이를테면 업무연속성계획(BCP) 시나리오의 다양화나 사이버 보험 등이 병행되어야 함을
시사한다. 이중채널 감쇠가 분포 수준에서 발현되는 전환점과 그 정량적 크기(꼬리
영역에서의 관측성 효과 감쇠 비율, 조건부 위험가치 격차 등)는 본 연구의 실증 결과에서
분할 회귀를 통해 확인되며, 이는 \secref{sec:ch5_dual_channel}에서 상세히 분석한다.
  % 2.5 시그모이드 함수와 바운딩


%% ===================================================================
\section{선행연구 검토 및 연구의 차별점}
\label{sec:lit-gap}
%% ===================================================================

앞에서 검토한 BCM, 정적 RTO의 한계, 관측성, 불확실성, 시그모이드 바운딩의 이론적 논의를 종합하고 선행연구의 흐름을 정리하여 본 연구가 수행하고자 하는 차별점을 [표 2-8]에 주제별로 정리하였다.

%% -------------------------------------------------------------------
%% -------------------------------------------------------------------

\begin{table}[htbp]
\centering
\caption{주제별 선행연구 요약}
\label{tab:prior-research}
\small
\begin{tabularx}{\textwidth}{@{}>{\raggedright\arraybackslash}p{1.8cm}
  >{\raggedright\arraybackslash}p{3.0cm}
  >{\raggedright\arraybackslash}X
  >{\raggedright\arraybackslash}p{4.0cm}@{}}
\toprule
\textbf{주제} & \textbf{주요 문헌} & \textbf{핵심 내용} & \textbf{한계} \\
\midrule
BCM/BCMS 표준 &
\citet{herbane2010disaster}, ISO 22301(2019) &
PDCA 기반 BCMS 프레임워크, BCM 발전 3단계 분류 &
RTO를 정적 고정값으로 설정, 동적 조정 메커니즘 부재 \\
\addlinespace
BIA 및 RTO &
\citet{snedaker2013bcdr}, \citet{hiles2010bcm} &
BIA 방법론 체계화, MTPD-RTO 관계 정립 &
RTO 동적 특성 미논의, 관측성 미반영 \\
\addlinespace
\m{관측성(제어 이론)} &
\citet{kalman1960control}, \citet{avizienis2004basic} &
LTI 시스템 관측성 형식화, 의존성 분류 체계 &
BCM 영역 적용 부재, RTO와의 수리적 연결 미시도 \\
\addlinespace
\m{관측성(IT 운영)} &
\citet{beyer2016sre}, \citet{majors2022observability} &
SRE 관측성 실무, 세 가지 핵심 요소(three pillars) 체계 &
정량적 모델화 미수행, MTTD/MTTR 관계만 정성적 \\
\addlinespace
가우시안 불확실성 &
\citet{jaynes2003probability}, \citet{gelman2013bayesian} &
CLT 기반 불확실성 정량화, 베이지안 추론 체계 &
극단값 과소평가, 파레토 현상 미포착 \\
\addlinespace
파레토 분포/EVT &
\citet{taleb2007blackswan}, \citet{embrechts1997extremal}, \citet{newman2005powerlaws} &
블랙 스완 개념, EVT 수학적 기초, 멱법칙 보편성 &
금융/보험 중심 적용, BCM/RTO 영역 미적용 \\
\addlinespace
국내 BCM &
\citet{cheung2019quickhit}, \citet{jang2021bcms}, \citet{jung2023bcms} &
BIA Quick-Hit 프레임워크, BCMS 성과 실증, BCM 평가지표 개발 &
정적 RTO 전제, 동적 조정 모델 미제시 \\
\addlinespace
DRTO, 시그모이드 바운딩 &
\citet{lee2026drto} &
관측성 기반 DRTO에 시그모이드 바운딩 도입, 곡률 매개변수 $\kappa^*$ 도출 &
관측성 단일 채널 한정, 불확실성 채널 미포함 \\
\bottomrule
\end{tabularx}
\end{table}

[표 2-8]에서 확인되는 선행연구의 공통적 한계는 \chg{다음 세 가지로 요약되며 2.7(DRTO 모델 구조)에서 상세히 다룬다. 첫째는 관측성을 통합한 RTO 모델의 부재로서 부분적인 단일 채널
모델만 존재하고~\citep{lee2026drto}, 둘째는 불확실성이 정량화된 동적 RTO 모델의 부재이며,
셋째는 극단 시나리오에 대한 파레토 분석의 부재이다.} 본 연구의 $\DRTO$ 프레임워크는 이
\chg{세 가지 한계}를 통합적으로 해소하여 선행연구와 차별화하는 것을 목표로 한다.


%% -------------------------------------------------------------------
%% -------------------------------------------------------------------

\p{[}표~\ref{tab:prior-research}\p{]}의 선행연구를 종합하면, 기존 BCM/RTO 연구와
본 연구의 $\DRTO$ 프레임워크 사이에는 다음 세 가지 핵심적 차별점이 존재한다.

\chg{첫째, 관측성을 통합한 RTO 모델의 부재이다.} 기존 BCM 연구에서 RTO는 BIA를 통해 고정적으로 결정되며, 시스템의 관측성 수준이 복구 시간에 미치는 영향은 정량적으로 모델링되지 않았다. 제어 이론\citep{kalman1960control}과 SRE\citep{beyer2016sre, majors2022observability}에서 관측성의 중요성이 \m{인정되며}, \citet{lee2026drto}은 관측성 단일 채널을 시그모이드 바운딩으로 RTO에 연결하는 수리적 프레임워크를 처음 제시하였으나, 불확실성 채널과 결합한 다채널 통합 모델은 여전히 부재하였다.

\chg{둘째, 불확실성이 정량화된 동적 RTO($\DRTO$) 모델의 부재이다.} 기존 RTO는 결정론적(deterministic) 단일값으로 설정되어 복구 과정에서 발생하는 불확실성을 반영하지 못하였다. 확률론적 사고방식\citep{gelman2013bayesian, jaynes2003probability}이 위험관리 분야에 일부 \m{적용되나}, RTO 자체를 확률 분포로 모델링하여 불확실성에 따른 동적 조정을 수행하는 연구는 부재하였다.

\chg{셋째, 극단적 시나리오에 대한 파레토 분석의 부재이다.} \m{극단값이론(EVT)}과 파레토 분포는 금융 위험관리\citep{embrechts1997extremal, mcneil2015risk}에서 활발히 \m{활용되나}, BCM 영역에서 극단적 복구 지연 시나리오를 파레토 분포로 모델링한 연구는 찾아보기 어렵다. \citet{taleb2007blackswan}이 지적한 극단 사건의 과소평가 문제가 RTO 설정에서도 동일하게 \m{발생함에도} 불구하고, 이에 대한 체계적 분석 및 대응 방안은 선행 연구에서 다루어지지 않았다.

[표 2-9]는 이상의 논의를 종합하여 기존 연구의 접근 방식과 본 연구의 DRTO 프레임워크가 제시하는 해결 방안을 대비하여 비교한 것이다.

\begin{table}[htbp]
\centering
\caption{선행 연구의 한계와 본 연구의 접근 비교}
\label{tab:research-gap}
\small
\begin{tabularx}{\textwidth}{@{}>{\raggedright\arraybackslash}p{2.6cm}
  >{\raggedright\arraybackslash}X
  >{\raggedright\arraybackslash}X
  >{\raggedright\arraybackslash}p{2.3cm}@{}}
\toprule
\textbf{연구 차별점} & \textbf{기존 연구} & \textbf{본 \m{연구($\DRTO$)}} & \textbf{관련 조건} \\
\midrule
차별점 1:\newline 관측성 미반영 &
단일 채널 관측성 모델~\citep{lee2026drto}만 존재;\newline 다채널 관측성-불확실성\newline 통합 부재 &
$O_{\text{raw}} \in [0,1]$ 관측성 원시값을\newline 개별 시그모이드 바운딩의\newline 핵심 입력으로 통합 &
\p{C1 (관측성)} \\
\addlinespace
차별점 2:\newline 불확실성 미정량 &
결정론적 단일값 RTO;\newline 확률적 변동 미고려 &
가우시안 $N(\mu,\sigma^2)$ 및\newline 파레토 $\text{Pareto}(\alpha)$\newline 이중채널 모델링 &
\p{C2 (가우시안),\newline C3 (파레토)} \\
\addlinespace
차별점 3:\newline 극단 시나리오\newline 미분석 &
BCM 영역에서 EVT/\newline 파레토 미적용;\newline 꼬리 위험 미포착 &
EVT 기반 파레토 분석;\newline 분할 회귀(Core/Tail);\newline 이중채널 비선형 효과 정량화 &
\p{C3 tail (파레토\newline P85.8 이상)} \\
\addlinespace
바운딩 기법 &
클리핑(hard boundary)\newline 또는 미적용 &
개별 시그모이드 바운딩\newline $S(z) = 2/(1+e^{-z})-1$\newline $\kappa = 1.2$, 전 구간 미분 가능 &
전 조건 공통 \\
\bottomrule
\end{tabularx}
\end{table}

\chg{[표 2-9]의 관련 조건 열에 나타나는 조건 표기는 관측성과 불확실성의 활성화 조합에 따라 정의한 네 가지 실험 조건을 가리킨다. C0는 관측성과 불확실성을 모두 비활성화한 기준선이고, C1은 관측성만 활성화한 조건이며, C2는 C1에 가우시안 불확실성을 더한 조건, C3는 C1에 파레토 불확실성을 더한 조건이다.}

본 연구의 $\DRTO$ 프레임워크는 위의 세 가지 차별점을 통합적으로 구현하는 것을 목표로 한다.
구체적으로, 관측성을 핵심 입력 변수로 도입하고,
가우시안 및 파레토 분포를 통해 불확실성을 이중 채널로 정량화하며,
파레토 분포를 BCM 영역에 도입하여 극단 시나리오 분석의 공백을 채운다.
다음으로 이러한 이론적 기반 위에 구축되는 DRTO 프레임워크의 구체적 수리 모델과 가설 체계를 상세히 기술한다.
  % 2.6 선행연구 검토 및 연구의 차별점
          % 2.1--2.6 이론적 배경
%% ===================================================================
\section{DRTO 모델 구조}
\label{sec:drto-model}
%% ===================================================================

앞서 도출한 이론적 배경과 선행 연구의 한계를 토대로, 본 절부터는 동적
복구목표시간($\DRTO$) 프레임워크의 수학적 구조와 설계 원리를 체계적으로 기술한다.
이후 계속하여 구체적으로 (i) 시그모이드 정식화(2.7), (ii) 4단계 조건 체계 C0--C3(2.8),
(iii) 수학적 성질(2.9), (iv) 곡률 매개변수 $\kappa^{*}=1.2$ 선정 근거(2.10),
(v) \jrev{10개 연구가설} H1--\jrev{H10}(2.11)을 차례로 다룬다.

%% -------------------------------------------------------------------
\subsection{\chg{개별 시그모이드 바운딩 핵심 수식}}
\label{subsec:core-equation}
%% -------------------------------------------------------------------

$\DRTO$ 모델의 핵심 과제는 관측성과 불확실성이라는 두 입력 차원의 효과를
물리적으로 타당한 범위 내에서 결합하는 것이다. 복구목표시간은 음수가 될 수 없으며,
조직이 허용할 수 있는 최대허용중단기간(MTPD)을 초과해서도 안 된다.
이러한 물리적 제약 조건을 만족하면서도 입력 변수의 변화에 연속적이고
미분 가능한 응답을 제공하는 모델을 구축하는 것이 설계의 핵심 요구사항이다.

본 연구에서는 \textbf{개별 시그모이드 바운딩(individual sigmoid bounding)} 모델을 제안한다.
선행연구\citep{lee2026drto}에서 수립한 관측성 기반 모델(C0$\to$C1)과
최적 곡률 $\kappa^{*} = 1.2$를 기반으로,
본 연구에서는 불확실성 채널(C2, C3)을 추가하여 다채널 구조로 확장한다.
관측성과 불확실성의 효과를 각각 고유한 물리적 범위(span)에서
독립적으로 바운딩하여, 물리적 경계 보장, 연속 미분 가능성,
이론적 유도 가능성이라는 세 요구사항을 동시에 충족한다.

\chg{모델의 기저 함수는} 수정 시그모이드(modified sigmoid)로서,
표준 로지스틱 함수를 $(-1, 1)$ 구간으로 사상한다.

\begin{equation}
  S(z) = \frac{2}{1 + e^{-z}} - 1, \quad S: \mathbb{R} \to (-1,\; 1)
  \label{eq:sigmoid}
\end{equation}

\chg{관측성의 효과는} 시그모이드 함수를 통해 $(-R_0, 0]$ 범위로 바운딩된다\chg{. 이때 $0$은
달성 가능한 값이지만 $-R_0$에는 한없이 접근할 뿐 도달하지는 못한다}.
관측성은 $\DRTO$를 $R_0$에서 최대 $0$까지 감소시킬 수 있으므로,
물리적 하향 범위(span)는 $R_0$이다.
\chg{이와 대칭적으로 불확실성의 효과는} 시그모이드 함수를 통해 $[0, R_{\max} - R_0)$ 범위로 바운딩된다\chg{. $0$은
달성 가능하나 $R_{\max} - R_0$에는 도달하지 못하며, $R_0 = 6.0$h와 $R_{\max} = 24.0$h를
가정하면 그 범위는 $[0, 18.0)$h이다}.
불확실성은 $\DRTO$를 $R_0$에서 최대 $R_{\max}$(MTPD)까지 증가시킬 수 있으므로,
물리적 상향 범위(span)는 $R_{\max} - R_0$이다.

\medskip
두 채널의 바운딩 수식은 다음과 같이 대칭적 구조를 갖는다.
\begin{align}
  E_{O,\text{bnd}} &= -R_0 \cdot S(z_O),
    &\quad \text{span} &= R_0
    \label{eq:E_O_bnd} \\[4pt]
  E_{U,\text{bnd}} &= \;(R_{\max} - R_0) \cdot S(z_U),
    &\quad \text{span} &= R_{\max} - R_0
    \label{eq:E_U_bnd}
\end{align}
시그모이드 입력 인자 $z_O$, $z_U$는 각각
\begin{align}
  z_O &= \kappa \cdot k_O \cdot O_{\text{raw}}
    \label{eq:z_O} \\[4pt]
  z_U &= \frac{\kappa \cdot k_U \cdot R_0 \cdot U_{\text{raw}}}{R_{\max} - R_0},
    \quad k_U = 1 - k_O
    \label{eq:z_U}
\end{align}
이다. $z_O = \kappa \cdot k_O \cdot O_{\text{raw}} \geq 0$은
입력 변수의 정의역에 의해 항상 성립하며,
$U_{\text{raw}} \geq 0$일 때 $z_U \geq 0$이므로
$E_{O,\text{bnd}} \in (-R_0, 0]$,\;
$E_{U,\text{bnd}} \in [0, R_{\max} - R_0)$이다.
$N(3,1)$에서의 음수 $U_{\text{raw}}$ 처리에 대해서는
\ref{subsec:c2}항에서 후술한다.

\chg{두 입력 인자 $z_O$, $z_U$는 세 가지 원칙에 따라 설계된다. 첫째는 무차원
정규화이다. 관측성 점수와 불확실성 원시값은 단위와 척도가 서로 다르므로, 각각
가중치를 곱한 뒤 불확실성은 물리적 상향 범위 $R_{\max} - R_0$로 나누어 두 입력을
모두 무차원량으로 변환한다. 이로써 시그모이드가 비교 가능한 척도 위에서
작동한다. 둘째는 물리적 범위 분리이다. 관측성은 복구 시간을 기준선 아래로
($\mathrm{span} = R_0$), 불확실성은 기준선 위로($\mathrm{span} = R_{\max} - R_0$)
작용하도록 두 채널의 방향과 범위를 분리한다. 셋째는 상보적 가중이다. 관측성
가중치와 불확실성 가중치가 $k_O + k_U = 1$을 이루어, 한 채널의 비중이 커지면
다른 채널의 비중이 줄어드는 상보 관계를 갖는다.}

두 입력 인자는 표면적 형태는 다르나 동일한 구조 $z = \kappa \cdot k \cdot (R_0/\mathrm{span})
\cdot X_{\text{raw}}$를 공유한다. 관측성 채널은 $\mathrm{span} = R_0$이므로 정규화 인자가
$R_0/R_0 = 1$이 되어 수식에서 생략되고, 불확실성 채널은 $\mathrm{span} = R_{\max} - R_0$이므로
정규화 인자가 $R_0/(R_{\max}-R_0) = 1/3$로 나타난다. 두 원시 입력은 모두 무차원이며,
시간 단위(h)는 바운딩 수식의 span에 의해 부여된다.
\p{[}표~\ref{tab:zo-zu-symmetry}\p{]}은 두 채널의 항목별 대칭 구조를 정리한 것이다.

\begin{table}[htbp]
\centering
\caption{관측성 채널($z_O$)과 불확실성 채널($z_U$)의 대칭 구조 비교}
\label{tab:zo-zu-symmetry}
\small
\begin{tabular}{l c c}
\toprule
\textbf{구분} & \textbf{관측성 ($z_O$)} & \textbf{불확실성 ($z_U$)} \\
\midrule
원시 입력     & $O_{\text{raw}} \in [0,1]$ (무차원) & $U_{\text{raw}} \geq 0$ (무차원) \\
정규화 인자   & $R_0 / R_0 = 1$              & $R_0 / (R_{\max} - R_0) = 1/3$ \\
인자 수식     & $\kappa \cdot k_O \cdot (R_0/R_0) \cdot O_{\text{raw}}$
              & $\kappa \cdot k_U \cdot (R_0/(R_{\max}\!-\!R_0)) \cdot U_{\text{raw}}$ \\
가중치        & $k_O$                               & $k_U = 1 - k_O$ \\
span (바깥)   & $R_0$                               & $R_{\max} - R_0$ \\
효과 방향     & 하향 ($\DRTO$ 감소)                 & 상향 ($\DRTO$ 증가) \\
\bottomrule
\end{tabular}
\end{table}

\chg{최종 $\DRTO$는} 기준선 $R_0$에 개별 바운딩된 두 효과를 가산하여 산출한다.

\begin{equation}
  \DRTO = R_0 + E_{O,\text{bnd}} + E_{U,\text{bnd}}
  \label{eq:drto}
\end{equation}

두 바운딩 효과의 범위가 물리적으로 분리되어 있으므로,
$\DRTO$는 항상 다음의 경계 내에 위치한다.

\begin{equation}
  E_{O,\text{bnd}} \in (-R_0, 0], \quad
  E_{U,\text{bnd}} \in [0, R_{\max} - R_0)
  \quad \therefore\;\; \DRTO \in (0,\; R_{\max})
  \label{eq:bounds}
\end{equation}

이 경계는 $\kappa$ 값에 무관하게 항상 성립하며,
$R_{\max} = 4R_0 = 24.0$h로 설정할 경우 $\DRTO \in (0,\; 24.0)$h가 보장된다.
특히 C1 조건($E_{U,\text{bnd}} = 0$)에서는 $\DRTO \in (0, R_0]$이 되어
$R_{\max}$가 불필요하다. 이는 관측성만의 효과를 불확실성 상한과
독립적으로 분석할 수 있게 하는 개별 바운딩의 핵심 장점이다.


\chg{조건 간 비교를 위해} 무차원 정규화 비율 $\eta$를 정의한다.

\begin{equation}
  \eta = \frac{\DRTO}{R_0}
  \label{eq:eta}
\end{equation}

$\eta = 1$은 정적 기준선과 동일함을, $\eta < 1$은 관측성에 의한 복구 시간 단축을,
$\eta > 1$은 불확실성에 의한 복구 시간 연장을 의미한다.
$\DRTO \in (0, R_{\max})$이므로 $\eta$의 이론적 범위는
$(0,\; R_{\max}/R_0)$이며, $R_{\max} = 4R_0$ 가정 하 $(0,\; 4)$이다.


%% -------------------------------------------------------------------
\subsection{모델 구조 도식}
\label{subsec:model-diagram}
%% -------------------------------------------------------------------

\p{[}그림~\ref{fig:model-structure}\p{]}는 $\DRTO$ 개별 시그모이드 바운딩 모델의
전체 구조를 도식화한 것이다. 두 입력 차원(관측성, 불확실성)이
각각 독립적인 시그모이드 바운딩 경로를 통과한 후 기준선 $R_0$에
가산되어 최종 $\DRTO$를 산출하는 구조를 보여준다.
\chg{그림에서 관측성 경로는 상단에 파란색으로, 불확실성 경로는 하단에 주황색으로 표시되며, 두 경로는 각각 독립적인 span에서 시그모이드 바운딩을 거쳐 가산 노드에서 기준선 $R_0$와 합산된다. 곡률 매개변수 $\kappa = 1.2$는 두 경로에 공통으로 적용된다.}

\begin{figure}[htbp]
\centering
\resizebox{0.95\textwidth}{!}{%
\begin{tikzpicture}[
  inputbox/.style={draw, rounded corners=5pt, minimum width=26mm,
    minimum height=13mm, align=center, font=\small, line width=0.6pt},
  procbox/.style={draw, rounded corners=4pt, minimum width=36mm,
    minimum height=12mm, align=center, font=\small, line width=0.6pt},
  outbox/.style={draw, rounded corners=5pt, minimum width=40mm,
    minimum height=16mm, align=center, font=\small\bfseries,
    line width=0.8pt},
  sumbox/.style={draw, circle, minimum size=10mm, align=center,
    font=\normalsize\bfseries, line width=0.8pt},
  arr/.style={-{Stealth[length=2.5mm]}, thick},
  lbl/.style={font=\scriptsize, fill=white, inner sep=1.5pt}
]

% ── 관측성 경로 (상단, y=3.0) ──
\node[inputbox, fill=blue!10] (ko) at (0, 4.0) {\textbf{가중치} $k_O$\\$\in [0,1]$};
\node[inputbox, fill=blue!10] (oraw) at (0, 2.0) {\textbf{원시 관측성} $O_{\text{raw}}$\\$\in [0,1]$};
\node[procbox, fill=blue!5] (zo) at (5.0, 3.0) {$z_O = \kappa \cdot k_O \cdot O_{\text{raw}}$};
\node[procbox, fill=blue!12] (sobnd) at (10.5, 3.0)
  {$E_{O,\text{bnd}} = -R_0 \cdot S(z_O)$\\{\scriptsize span $= R_0$}};

% ── 기준선 (중앙, y=0) ──
\node[procbox, fill=gray!15] (r0) at (5.0, 0) {기준선 $R_0 = 6.0$h};

% ── 불확실성 경로 (하단, y=-3.0) ──
\node[inputbox, fill=orange!10] (ku) at (0, -1.5) {\textbf{가중치} $k_U$\\$= 1 - k_O$};
\node[inputbox, fill=orange!10] (unorm) at (0, -4.0) {\textbf{원시 불확실성} $U_{\text{raw}}$\\$\geq 0$};
\node[procbox, fill=orange!5] (zu) at (5.0, -3.0)
  {$z_U = \kappa \cdot k_U \cdot \frac{R_0 \cdot U_{\text{raw}}}{R_{\max} - R_0}$};
\node[procbox, fill=orange!12] (subnd) at (10.5, -3.0)
  {$E_{U,\text{bnd}} = (R_{\max}\!-\!R_0) \cdot S(z_U)$\\{\scriptsize span $= R_{\max} - R_0$}};

% ── 가산 노드 ──
\node[sumbox, fill=green!8] (sum) at (14.5, 0) {$+$};

% ── 출력 ──
\node[outbox, fill=green!8] (drto) at (19.0, 0)
  {$\DRTO = R_0 + E_{O} + E_{U}$\\$\eta = \DRTO / R_0 \in (0,\; 4)$};

% ── 화살표: 관측성 경로 ──
\draw[arr, blue!60] (ko.east) -- (zo.150);
\draw[arr, blue!60] (oraw.east) -- (zo.210);
\draw[arr, blue!70] (zo.east) -- (sobnd.west);

% ── 화살표: 불확실성 경로 ──
\draw[arr, orange!60] (ku.east) -- (zu.150);
\draw[arr, orange!60] (unorm.east) -- (zu.210);
\draw[arr, orange!70] (zu.east) -- (subnd.west);

% ── 화살표: 3경로 → 가산 노드 (진입점 분리, 겹침 방지) ──
\draw[arr, blue!70]   (sobnd.east) -- ++(0.8,0) |- (sum.130);
\draw[arr, black]  (r0.east) -- ++(6.5,0) -- (sum.west);
\draw[arr, orange!70] (subnd.east) -- ++(0.8,0) |- (sum.230);

% ── 화살표: 가산 노드 → 출력 ──
\draw[arr, black] (sum.east) -- (drto.west);

% ── 매개변수 κ (중앙, 두 경로 사이) ──
\node[draw, rounded corners=3pt, fill=red!5, inner sep=3pt,
  font=\scriptsize, text=black] at (5.0, 0.9)
  {$\kappa = 1.2$};

% ── κ → 두 z 노드 연결 (점선) ──
\draw[densely dashed, red!40, thin, -{Stealth[length=1.5mm]}]
  (5.0, 0.7) -- (zo.south);
\draw[densely dashed, red!40, thin, -{Stealth[length=1.5mm]}]
  (5.0, -0.7) -- (5.0, -2.2);

% ── 범위 라벨 ──
\node[font=\scriptsize, text=black] at (10.5, 4.5)
  {$E_{O,\text{bnd}} \in (-R_0,\; 0]$};
\node[font=\scriptsize, text=black] at (10.5, -4.5)
  {$E_{U,\text{bnd}} \in [0,\; R_{\max}-R_0)$};

% ── 물리적 경계 ──
\node[draw, dashed, rounded corners=3pt, inner sep=5pt,
  font=\scriptsize, text=black] at (19.0, -2.5)
  {물리적 경계: $0 < \DRTO < R_{\max} = 24.0$h};

\end{tikzpicture}%
}% end resizebox
\caption{$\DRTO$ 개별 시그모이드 바운딩 모델 구조}
\label{fig:model-structure}
\end{figure}


%% -------------------------------------------------------------------
\subsection{변수 및 매개변수 정의}
\label{subsec:parameters}
%% -------------------------------------------------------------------

\p{[}표~\ref{tab:parameters}\p{]}은 $\DRTO$ 모델의 전체 변수와 매개변수를 정리한 것이다.
각 매개변수의 설정 근거는 연구자의 논문 \citet{lee2026drto}, 국제 표준, 또는 수학적 제약 조건에
기반한다.

\begin{table}[htbp]
\centering
\caption{DRTO 모델 변수 및 매개변수 정의}
\label{tab:parameters}
\small
\begin{tabularx}{\textwidth}{l l c c X}
\toprule
\textbf{기호} & \textbf{명칭} & \textbf{값/범위} & \textbf{단위} & \textbf{설정 근거} \\
\midrule
$R_0$ & 기준선 RTO & 6.0 & h & \nrev{ISO 22301:2019} BIA 기반 표준 설정값 \\
$R_{\max}$ & DRTO 상한 (MTPD) & 24.0 & h & $R_{\max} = 4 \cdot R_0$ (MTPD 기반 상한) \\
$\kappa$ & 곡률 매개변수 & 1.2 & --- & 다기준 최적화($\kappa^{*}$)~\citep{lee2026drto} \\
$S(z)$ & 수정 시그모이드 & $(-1, 1)$ & --- & $S(z) = 2/(1+e^{-z})-1$ \\
$k_O$ & 관측성 가중치 & $[0, 1]$ & --- & $\mathrm{Uniform}(0, 1)$으로 시뮬레이션 \\
$O_{\text{raw}}$ & 관측성 원시값 & $[0, 1]$ & --- & $\mathrm{Uniform}(0, 1)$으로 시뮬레이션 \\
$k_U$ & 불확실성 가중치 & $[0, 1]$ & --- & $k_U = 1 - k_O$ (상보적) \\
$z_O$ & 관측성 인자 & $\geq 0$ & --- & $\kappa \cdot k_O \cdot O_{\text{raw}}$ \\
$z_U$ & 불확실성 인자 & $\geq 0$ & --- & $\kappa \cdot k_U \cdot R_0 \cdot U_{\text{raw}} / (R_{\max} - R_0)$ \\
$U_{\text{raw}}^{(G)}$ & Gaussian 불확실성 & $N(3, 1)$ & --- & 평균 $\mu=3$, 표준편차 $\sigma=1$ \\
$U_{\text{raw}}^{(P)}$ & Pareto 불확실성 & $6 \cdot \mathrm{Pareto}(3)$ & --- & $\alpha=3$, $x_m=1$, 스케일링($\times 6$) \\
$E_{O,\text{bnd}}$ & 관측성 바운딩 조정량 & $(-R_0, 0]$ & h & span $= R_0 = 6.0$h \\
$E_{U,\text{bnd}}$ & 불확실성 바운딩 조정량 & $[0, R_{\max}\!-\!R_0)$ & h & span $= R_{\max} - R_0 = 18.0$h \\
$\eta$ & 효율비 & $(0, 4)$ & --- & $\DRTO / R_0$ \\
\bottomrule
\end{tabularx}
\end{table}


%% -------------------------------------------------------------------
\subsection{물리적 경계 보장 증명}
\label{subsec:boundary-proof}
%% -------------------------------------------------------------------

\phantomsection\label{prop:boundary}%
임의의 $\kappa > 0$, $k_O \in [0, 1]$, $O_{\text{raw}} \in [0, 1]$,
$z_U \geq 0$에 대해 $0 < \DRTO < R_{\max}$가 성립한다.(명제 2.1)

이 명제의 증명은 다음과 같다.
$z_O = \kappa \cdot k_O \cdot O_{\text{raw}} \geq 0$이므로
$S(z_O) \in [0, 1)$이다. 따라서
$E_{O,\text{bnd}} = -R_0 \cdot S(z_O) \in (-R_0, 0]$이다.
$z_U \geq 0$이므로
$S(z_U) \in [0, 1)$이다. 따라서
$E_{U,\text{bnd}} = (R_{\max} - R_0) \cdot S(z_U) \in [0, R_{\max} - R_0)$이다.
이 두 부등식을 합산하면 $\DRTO = R_0 + E_{O,\text{bnd}} + E_{U,\text{bnd}} \in (0,\; R_{\max})$이므로
$0 < \DRTO < R_{\max}$가 $\kappa$에 무관하게 성립한다. (명제 증명 완료)

위 증명은 $z_U \geq 0$ (즉, $U_{\text{raw}} \geq 0$) 전제 하에 성립한다.
C2 조건에서 $N(3,1)$의 음수 표본($0.135\%$)은 $z_U < 0$을 유발하여
$E_{U,\text{bnd}} < 0$이 되나,
$|E_{O,\text{bnd}}| < R_0$이고 $|E_{U,\text{bnd}}| < R_{\max} - R_0$이므로
$\DRTO > R_0 - R_0 - (R_{\max} - R_0) = 2R_0 - R_{\max}$이다.
$R_{\max} = 4R_0$일 때 이론적 하한은 $-2R_0$이나,
실제 시뮬레이션에서 $\DRTO$ 최솟값은 $2.861$h로 $0$을 충분히 상회한다.
즉 이론 하한이 음수임에도 음수 표본 출현 확률이 $0.135\%$로 작아 실질적 영향은 무시 가능하다.

\medskip
\chg{\noindent 이 경계 보장은 세 가지 실무적 의의를 갖는다. 첫째, 양수
보장($\DRTO > 0$)으로 $\DRTO$가 음수가 되는 물리적으로 무의미한 상황을
원천적으로 방지한다. 둘째, 상한 보장($\DRTO < R_{\max}$)으로
\nrev{ISO 22301:2019}의 MTPD 제약을 수학적으로 충족한다.
셋째, 클리핑(hard bounding)과 달리 경계 근처에서 미분 가능하여 매개변수 변화에
연속적으로 응답한다.}


%% ===================================================================
\section{4개 실험 조건 (C0--C3)}
\label{sec:conditions}
%% ===================================================================

$\DRTO$ 모델의 검증을 위해 관측성과 불확실성의 활성화 조합에 따라
네 개의 실험 조건(C0--C3)을 체계적으로 정의한다.
이 중 C0(기준선)과 C1(관측성 활성화)은 \citet{lee2026drto}에서
\chg{수립되고 검증된} 조건이며, C2(가우시안 불확실성)와 C3(파레토 불확실성)는
본 연구에서 새롭게 도입한 조건이다.
조건 설계의 핵심 원리는 \textbf{점진적 활성화(progressive activation)}로서,
C0에서 C3까지 한 번에 하나의 요소를 추가하여 각 요소의 개별 효과와
결합 효과를 분리 분석할 수 있도록 한다.

\p{[}표~\ref{tab:conditions}\p{]}는 각 조건의 구성과 특성을 비교한 것이다.

\begin{table}[htbp]
\centering
\caption{4개 실험 조건(C0--C3) 정의 및 시뮬레이션 결과 요약}
\label{tab:conditions}
\small
\begin{tabularx}{\textwidth}{c l c c l c X}
\toprule
\textbf{조건} & \textbf{명칭} & \textbf{관측성} & \textbf{불확실성}
  & \textbf{$U_{\text{raw}}$ 분포} & \textbf{평균 $\eta$} & \textbf{수식} \\
\midrule
\textbf{C0} & 정적 기준선 & OFF & OFF & --- & 1.000 &
  $\DRTO = R_0$ \\
\textbf{C1} & 관측성 전용 & ON & OFF & --- & 0.853 &
  $\DRTO = R_0 + E_{O,\text{bnd}}$ \\
\textbf{C2} & Gaussian 불확실성 & ON & ON & $N(3, 1)$ & 1.682 &
  $\DRTO = R_0 + E_{O,\text{bnd}} + E_{U,\text{bnd}}^{(G)}$ \\
\textbf{C3} & Pareto 불확실성 & ON & ON & $6 \cdot \mathrm{Pa}(3)$ & 1.514 &
  $\DRTO = R_0 + E_{O,\text{bnd}} + E_{U,\text{bnd}}^{(P)}$ \\
\bottomrule
\end{tabularx}
\end{table}


\p{[}그림~\ref{fig:conditions-progression}\p{]}은 C0에서 C3까지의 조건 진행과
각 단계에서 활성화되는 요소를 도식화한 것이다.
\chg{정적 조건 C0에서 관측성을 도입한 C1로 이어지는 공통 경로를 거친 뒤, 가우시안 불확실성의 C2와 파레토 불확실성의 C3가 C1에서 병렬로 분기하는 구조이다.}

\begin{figure}[htbp]
\centering
\resizebox{0.95\textwidth}{!}{%
\begin{tikzpicture}[
  cond/.style={draw, rounded corners=6pt, minimum width=30mm,
    minimum height=20mm, align=center, font=\small, line width=0.7pt},
  arr/.style={-{Stealth[length=3mm]}, thick},
  arrlbl/.style={font=\scriptsize, fill=white, inner sep=2pt,
    rounded corners=1pt},
  comp/.style={font=\scriptsize, text=black, align=center},
  efflbl/.style={font=\scriptsize, align=center, text=black}
]

% 조건 노드
\node[cond, fill=gray!15] (c0) at (0, 0)
  {\textbf{C0}\\정적 기준선\\$\eta = 1.000$\\$\DRTO = 6.0$h};
\node[cond, fill=blue!12] (c1) at (5.5, 0)
  {\textbf{C1}\\관측 전용\\$\bar{\eta} = 0.853$\\$\overline{\DRTO} = 5.12$h};
\node[cond, fill=green!12] (c2) at (11, 0)
  {\textbf{C2}\\Gaussian 불확실성\\$\bar{\eta} = 1.682$\\$\overline{\DRTO} = 10.09$h};
\node[cond, fill=red!12] (c3) at (16.5, 0)
  {\textbf{C3}\\Pareto 불확실성\\$\bar{\eta} = 1.514$\\$\overline{\DRTO} = 9.09$h};

% 화살표 및 라벨 (상단: 전이 조건, 중앙: Cohen's d)
\draw[arr] (c0) -- (c1);
\node[arrlbl, above=3pt] at ($(c0)!0.5!(c1)$) {$+O_{\text{raw}},\;k_O$};
\node[efflbl] at ($(c0)!0.5!(c1) + (0, 1.3)$)
  {Cohen's $d = -1.180$\\(C1 vs C0)};

\draw[arr] (c1) -- (c2);
\node[arrlbl, above=3pt] at ($(c1)!0.5!(c2)$) {$+U \sim N(3,1)$};
\node[efflbl] at ($(c1)!0.5!(c2) + (0, 1.3)$)
  {Cohen's $d = +1.871$\\(C2 vs C1)};

%% 분기 B: C1 → C3 (ㄷ자 직각 경로: 아래→오른쪽→위)
\draw[arr] (c1.south) -- ++(0,-0.6)
          -- ($(c3.south) + (0,-0.6)$) -- (c3.south);
\node[arrlbl, above=2pt] at ($(c1.south)!0.5!(c3.south) + (0,-0.6)$)
  {$+U \sim 6\!\cdot\!\mathrm{Pa}(3)$};
\node[efflbl] at ($(c2)!0.5!(c3) + (0, 1.3)$)
  {CVaR$_{99}$ 증가\\(C3 vs C2)};

% 하단 구성 요소
\node[comp, below=5mm of c0] {관측성: OFF\\불확실성: OFF};
\node[comp, below=5mm of c1] {관측성: ON\\불확실성: OFF};
\node[comp, below=5mm of c2] {관측성: ON\\불확실성: 가우시안(경꼬리)};
\node[comp, below=5mm of c3] {관측성: ON\\불확실성: 파레토(중꼬리)};

\end{tikzpicture}%
}% end resizebox
\caption{4개 실험 조건의 점진적 활성화}
\label{fig:conditions-progression}
\end{figure}


%% -------------------------------------------------------------------
\subsection{\chg{정적 기준선 C0}}
\label{subsec:c0}
%% -------------------------------------------------------------------

C0 조건은 관측성과 불확실성이 모두 비활성화된 상태로서,
기존 \nrev{ISO 22301:2019} 기반 정적 RTO 체계를 재현한다.
$O_{\text{raw}} = 0$, $U_{\text{raw}} = 0$이므로 $z_O = z_U = 0$이 되어
$S(0) = 0$이 성립하며, 따라서 $\DRTO = R_0 = 6.0$h, $\eta = 1.000$이다.
C0는 후속 조건들의 비교 기준선(baseline)으로 기능하며,
모든 동적 효과의 출발점을 제공한다.


%% -------------------------------------------------------------------
\subsection{\chg{C1의 관측성 기반 동적 조정}}
\label{subsec:c1}
%% -------------------------------------------------------------------

C1 조건은 관측성만 활성화하고 불확실성은 비활성화한 상태이다.
$U_{\text{raw}} = 0$이므로 $E_{U,\text{bnd}} = 0$이 되어
$\DRTO = R_0 - R_0 \cdot S(\kappa \cdot k_O \cdot O_{\text{raw}})$이다.
$S(\cdot) \geq 0$이므로 $\DRTO \leq R_0$가 보장되며,
이는 관측성 활성화가 복구 시간을 단축하는 방향으로만 작용함을 의미한다.
개별 바운딩의 핵심 장점으로, C1에서는 $R_{\max}$가 수식에 등장하지 않아
관측성만의 효과를 불확실성 상한과 독립적으로 분석할 수 있다.

$k_O$와 $O_{\text{raw}}$가 모두 최대값(1.0)일 때
$z_O = 1.2$가 되어
$S(1.2) \approx 0.537$이므로
$\DRTO \approx 2.78$h가 된다.
이는 완전 관측 상태에서 약 $53.7\%$의 복구 시간 단축이 가능함을 나타낸다.

\chg{이처럼 C1에서는 관측성 활성화가 복구 시간을 단축하는 방향으로만 작용하며, 그 평균적 단축 효과와 효과 크기가 기준선(C0) 대비 통계적으로 유의한가는 본 장이 제시한 이론적 예측에 해당한다. 표본 기반 실증값(평균 효율비, 단축률, Cohen's $d$)은 제4장에서 정량적으로 검증한다(\secref{sec:result_c0c1} 참조).}


%% -------------------------------------------------------------------
\subsection{\chg{C2의 가우시안 불확실성 결합}}
\label{subsec:c2}
%% -------------------------------------------------------------------

C2 조건은 관측성에 더하여 가우시안(Gaussian) 불확실성을 결합한 상태이다.
불확실성 원시값은 $U_{\text{raw}}^{(G)} \sim N(3, 1)$으로 생성되며,
$z_U = \kappa \cdot k_U \cdot R_0 \cdot U_{\text{raw}}^{(G)} / (R_{\max} - R_0)$으로 변환한 후
수식~(\ref{eq:E_U_bnd})에 따라 불확실성 효과 $E_{U,\text{bnd}}$를 산출한다.

\chg{한편 가우시안 분포는 음수 표본을 낼 수 있어 그 처리를 짚어 둔다.} 가우시안 분포의 지지(support)는 $(-\infty, +\infty)$이므로
$U_{\text{raw}}^{(G)} < 0$인 표본이 이론적으로 발생할 수 있다.
그러나 $N(3,1)$에서 $\Pr(U < 0) = \Phi(-3) \approx 0.135\%$로서
$n = 10{,}000$ 표본 중 13건에 불과하다.
해당 표본에서는 $z_U < 0$이 되어 $E_{U,\text{bnd}} < 0$이 산출되나,
실측 범위는 $[-1.13,\; -0.01]$h에 그치며
$\DRTO$ 최솟값이 $2.861$h로 물리적 하한($0$)을 충분히 상회한다.
따라서 절단정규분포(truncated normal)나 클리핑(clipping) 등의
인위적 분포 변형을 적용하지 않는다.
\chg{이는 세 가지 근거에 의한다. 첫째, 빈도 측면에서 음수 표본은 $0.13\%$(실측
13/10{,}000, 이론 $0.135\%$)에 불과하여 통계적 결론에 무시할 수 있는 영향만 미친다.
둘째, 경계 측면에서 해당 표본에서도 $\DRTO = 2.861$h로 0보다 충분히 크다. 셋째, 분포
보존 측면에서 절단을 적용하면 $N(3,1)$의 확률 구조가 왜곡되어 C2와 C3 사이의 공정한
비교(동일 기대값 $3.0$ 설계)가 훼손된다.}

가우시안 분포는 대칭적이고 경꼬리(light-tail) 특성을 가지므로,
C2 조건은 일상적 운영 변동, \chg{곧 인력 가용성의 일시적 변화와
장비 성능의 통상적 편차, 절차 수행 시간의 자연적 변이를} 반영한다.

\chg{C2에서는 이러한 가우시안 불확실성의 결합이 복구 시간을 연장하는 방향으로 작용하리라 예측되며, 그 평균적 연장 크기와 C1 대비 효과 크기는 제4장에서 정량적으로 검증한다(\secref{sec:result_c2} 참조).}


%% -------------------------------------------------------------------
\subsection{\chg{C3의 파레토 Heavy-tail 결합}}
\label{subsec:c3}
%% -------------------------------------------------------------------

C3 조건은 C1에서 분기하여 가우시안(C2) 대신 파레토(중꼬리, heavy-tail) 분포로
불확실성을 모델링한 변종이다. 불확실성 원시값은
$U_{\text{raw}}^{(P)} \sim 6 \cdot \mathrm{Pareto}(\alpha = 3)$으로
생성되며, $z_U = \kappa \cdot k_U \cdot R_0 \cdot U_{\text{raw}}^{(P)} / (R_{\max} - R_0)$으로 변환된다.
$\alpha = 3$은 평균은 유한하나 분산이 가우시안($\sigma^2=1$) 대비 약 27배인 분포를 제공한다.

C2와 C3의 불확실성 분포는 모집단 평균($\mu = 3.0$)이 동일하도록 설계되어
공정한 비교가 가능하다. 두 분포의 차이는 분산과 극단 백분위에서 발현된다.
가우시안의 99백분위가 약 5.33인 반면, 파레토의 99백분위는 약 21.85로서
약 4.1배의 차이를 보인다.

\chg{파레토 분포는 중꼬리 특성을 가지므로, C3는 평균 수준에서는 C2와 유사하더라도 극단 백분위로 갈수록 $\DRTO$가 C2를 크게 상회하리라 예측된다. 두 분포의 누적분포함수가 교차하는 지점이 존재하며, 본 연구에서는 이를 ``CDF 교차점(crossover)''으로 명명한다. 교차점을 경계로 핵심 영역(Core)과 꼬리 영역(Tail)을 분할하면, 꼬리 영역에서는 시그모이드 포화로 인해 관측성 가중치 $k_O$의 회귀 계수가 감쇠하는 이중채널 비선형 상호작용(dual-channel nonlinear interaction)이 나타나리라 예측된다. 교차점의 구체적 위치와 조건별 평균 효율비, 그리고 감쇠의 실증값은 \jrev{4.10절}(이중채널)에서 상세히 분석한다.}
           % 2.7--2.8 DRTO 모델·실험 조건


%% ===================================================================
\section{시그모이드 바운딩의 수학적 성질}
\label{sec:sigmoid-properties}
%% ===================================================================

본 절에서는 $\DRTO$ 모델의 기저 함수인 수정 시그모이드 $S(z)$의
수학적 성질을 체계적으로 분석한다.

%% -------------------------------------------------------------------
\subsection{기본 해석학적 성질}
\label{subsec:basic-properties}
%% -------------------------------------------------------------------

수정 시그모이드 함수 $S(z) = 2/(1+e^{-z}) - 1$은 다음의 성질을 갖는다.
\chg{앞서 2.5(시그모이드 함수와 바운딩)에서 바운딩 기법의 관점에서 개관한 기본 성질을, 2.9(시그모이드 바운딩의 수학적 성질)에서는 DRTO 바운딩 함수로서 2계 도함수와 변곡점, 곡률 거동까지 포함하여 엄밀하게 정리한다.}

\chg{첫째, 전역 미분 가능성($C^{\infty}$) 측면에서 $S(z)$는 지수 함수의 합성으로 구성되므로
전체 정의역에서 무한 번 미분 가능하며, 1차 도함수가 $S'(z) = 2e^{-z}/(1+e^{-z})^2 > 0$
($\forall z \in \mathbb{R}$)이므로 전 구간에서 단조증가한다. 둘째, 단조증가성에 의해 입력의
순서가 출력에 보존되어 $z_1 < z_2$이면 $S(z_1) < S(z_2)$가 성립하며, 이는 관측성이 높을수록
$\DRTO$가 더 감소하고 불확실성이 클수록 더 증가하는 물리적 직관과 일치한다. 셋째, 경계
보장 측면에서 $z \to +\infty$일 때 $S(z) \to +1$, $z \to -\infty$일 때 $S(z) \to -1$이므로
$S: \mathbb{R} \to (-1, 1)$이고, 점근선에 도달하되 등호가 성립하지 않으므로 $\DRTO$가 정확히
$0$ 또는 $R_{\max}$에 도달하는 것은 불가능하다. 넷째, 원점 대칭성(홀함수)로서
$S(-z) = -S(z)$이고 $S(0) = 0$이어서 효과가 없을 때 기준선을 유지하는 중립 조건을 보장하며,
$z_O = z_U = 0$일 때 $E_{O,\text{bnd}} = E_{U,\text{bnd}} = 0$이 되어 $\DRTO = R_0$로 기준선이
유지된다. 다섯째, 2계 도함수가 $S''(z) = -2 e^{-z} (1 - e^{-z})/(1 + e^{-z})^3$로서 $z = 0$에서
부호가 바뀌는 변곡점을 가지며, 이 변곡점이 ``효과 없음'' 지점($S(0)=0$)과 일치하므로 작은
입력에서는 함수 거동이 선형에 가깝다가 $|z|$가 커질수록 비선형성이 강화되는 구조를 갖는다.}


%% -------------------------------------------------------------------
\subsection{작은 $\kappa \cdot z$ 값 영역에서의 선형 근사}
\label{subsec:linear-approx}
%% -------------------------------------------------------------------

$|z|$가 충분히 작을 때, $S(z)$는 \m{식~(\ref{eq:linear-approx})에 나타낸 바와 같이} 1차 테일러 전개로 근사된다.

\begin{equation}
  S(z) \approx S'(0) \cdot z = \frac{1}{2} z, \quad |z| \ll 1
  \label{eq:linear-approx}
\end{equation}

이 근사는 $|z| \leq 0.5$에서 약 $2\%$ 이하의 오차로 성립하며,
$|z| \leq 0.8$에서는 약 $5\%$ 이하의 오차를 갖는다.

이 근사를 C1 수식에 적용하면
$E_{O,\text{bnd}} \approx -R_0 \cdot \frac{1}{2} \cdot \kappa \cdot k_O \cdot O_{\text{raw}}$
$= -\frac{\kappa}{2} \cdot R_0 \cdot k_O \cdot O_{\text{raw}}$가 되어
선형 모델과의 대응 관계가 명확해진다.

$\kappa = 1.2$일 때 $z_O = \kappa \cdot k_O \cdot O_{\text{raw}} \in [0, 1.2]$이다.
$z_O$의 상한($k_O = O_{\text{raw}} = 1$일 때 $z_O = 1.2$)에서 $S(1.2) = 0.537$이고
선형 근사는 $\frac{1}{2} \times 1.2 = 0.600$으로서, 근사 오차는 약 $10.5\%$이다.
따라서 $\kappa = 1.2$는 선형 영역에서 비선형 전이 영역으로 진입하는 경계에 위치하며,
시그모이드의 비선형 압축 효과(선형 근사 대비 출력값 감소)가 의미 있게 발현되면서도
포화에는 도달하지 않는 균형 영역에 해당한다.

\p{[}그림~\ref{fig:sigmoid-curve}\p{]}은 $\kappa = 1.2$에서의 시그모이드 함수
$S(z)$와 선형 근사 $z/2$를 비교한 것이다.
\chg{그림에서 관측성 인수 $z_O$의 범위 $[0, 1.2]$가 시그모이드의 선형 영역에서 비선형 전이 영역으로 확장됨을 볼 수 있으며, $S(1.2) = 0.537$로서 선형 근사값 $0.600$ 대비 약 $10.5\%$의 압축이 발생한다.}

\begin{figure}[htbp]
\centering
\begin{tikzpicture}
\begin{axis}[
  width=0.82\textwidth,
  height=0.50\textwidth,
  xlabel={$z$},
  ylabel={$S(z)$},
  xmin=-4, xmax=4,
  ymin=-1.15, ymax=1.15,
  grid=major,
  grid style={black},
  legend pos=north west,
  legend style={font=\small, rounded corners=2pt},
  axis lines=middle,
  every axis x label/.style={at={(axis description cs:1.02,0.5)},
    anchor=west, font=\small},
  every axis y label/.style={at={(axis description cs:0.5,1.04)},
    anchor=south, font=\small},
  clip=false
]
  % 시그모이드 곡선
  \addplot[thick, black, domain=-4:4, samples=200]
    {2/(1+exp(-x)) - 1};
  \addlegendentry{$S(z) = 2/(1+e^{-z})-1$}

  % 선형 근사
  \addplot[thick, dashed, blue!70, domain=-2.2:2.2, samples=100]
    {0.5*x};
  \addlegendentry{선형 근사 $z/2$}

  % 경계선
  \addplot[thin, dashed, black] coordinates {(-4, 1) (4, 1)};
  \addplot[thin, dashed, black] coordinates {(-4, -1) (4, -1)};

  % κ·k_O·O_raw 범위 표시
  \draw[thick, red!70, {Stealth[length=2mm]}-{Stealth[length=2mm]}]
    (axis cs:0, -1.08) -- (axis cs:1.2, -1.08);
  \node[font=\scriptsize, text=black, anchor=north] at (axis cs:0.6, -1.08)
    {$\kappa k_O O_{\text{raw}} \in [0, 1.2]$};

  % S(1.2) 표시
  \draw[thin, dotted, red!50] (axis cs:1.2, 0) -- (axis cs:1.2, 0.537);
  \node[font=\tiny, text=black, anchor=west] at (axis cs:1.25, 0.537)
    {$S(1.2)=0.537$};
  \filldraw[red!70] (axis cs:1.2, 0.537) circle (1.5pt);

  % 경계 라벨
  \node[font=\scriptsize, text=black, anchor=west] at (axis cs:4.1, 1) {$+1$};
  \node[font=\scriptsize, text=black, anchor=west] at (axis cs:4.1, -1) {$-1$};

\end{axis}
\end{tikzpicture}
\caption{수정 시그모이드 함수 $S(z)$와 선형 근사 $z/2$의 비교}
\label{fig:sigmoid-curve}
\end{figure}


%% ===================================================================
\section{최적 곡률 파라미터 $\kappa^{*} = 1.2$}
\label{sec:kappa-optimization}
%% ===================================================================

곡률 매개변수 $\kappa$는 시그모이드 바운딩 모델에서 입력 변수의 변화에 대한
$\DRTO$의 반응 강도를 결정하는 핵심 설계 변수이다.
$\kappa$가 너무 작으면 시그모이드가 선형 영역에 머물러
입력 변화에 대한 $\DRTO$의 반응이 미약하고(과소 반응),
$\kappa$가 너무 크면 적은 입력에서도 $|S|$가 1에 근접하여 다수 표본이
포화 영역에 진입함으로써 극단값 간의 구분력이 상실된다(조기 포화).
본 연구는 \citet{lee2026drto}의 다기준 최적화(multi-criteria
optimization) 방법을 적용하여 $\kappa^{*} = 1.2$를 도출하였다.


%% -------------------------------------------------------------------
\subsection{4개 최적화 기준}
\label{subsec:four-criteria}
%% -------------------------------------------------------------------

$\kappa$의 최적값 선정을 위해 다음의 4개 기준을 동시에 고려한다.

\chg{첫째는 효과 크기(Effect Size)로서 C1 대 C0의 Cohen's $d$ \m{절댓값}
$|d_{\text{C1vsC0}}|$이며, $\kappa$가 클수록 관측성의 효과가 증가하여 $|d|$가 커지므로
$|d| \geq 0.8$(대효과)을 기준으로 한다. 둘째는 포화율(Saturation Rate)로서 관측성
채널에서 $|S(z_O)| \geq \erf{0.99}{0.95}$인 표본의 비율이며, $\kappa$가 너무 크면 다수의 표본이 포화
영역에 진입하여 입력 변화에 대한 구분력이 상실되므로 포화율 $\leq 5\%$를 기준으로 한다.
셋째는 회귀 설명력($R^2$)으로서 C1에서의 2차 다항 회귀 결정계수이며, 시그모이드의
비선형성이 과도하면 2차 다항 근사의 정확도가 저하되므로 $R^2 \geq 0.95$를 기준으로 한다.
넷째는 RTO 단축률(Reduction Rate)로서 C1 평균 효율성 비율에 의한 단축률
$(1 - \bar{\eta}_{\text{C1}}) \times 100\%$이며, 실무적으로 의미 있는 단축 효과를 제공하는지를
평가하기 위해 단축률 $\geq 10\%$를 기준으로 한다.}


%% -------------------------------------------------------------------
\subsection{종합점수 $F(\kappa)$와 격자 탐색}
\label{subsec:composite-score}
%% -------------------------------------------------------------------

4개 기준을 종합하는 종합 점수(composite score) $F(\kappa)$를
\m{식~(\ref{eq:composite})에 나타낸 바와 같이} 곱함수(product function)로 정의한다.

\begin{equation}
  F(\kappa) = f_1(\kappa) \cdot f_2(\kappa) \cdot f_3(\kappa) \cdot f_4(\kappa)
  \label{eq:composite}
\end{equation}

여기서 $f_i(\kappa) \in [0, 1]$은 각 기준의 부분 점수이다.
각 부분 점수는 해당 기준에서 가능한 최대값으로 정규화되어 $f_i \in [0,1]$이며,
따라서 $F \in [0,1]$이고 $F = 1$은 모든 기준이 동시에 최적치를 달성한 상태를
의미한다.
곱 구조를 채택한 근거는 다음과 같다.
하나의 기준이라도 0이면 $F = 0$이 되어 모든 기준의 동시 충족을 강제하고,
기준 간 대체(trade-off)를 허용하지 않으며,
기준별 가중치 결정이 불필요하다.
$F(\kappa)$가 최대인 $\kappa$를 $\kappa^{*}$로 선정한다.

각 부분 점수 $f_i(\kappa)$의 구체적 함수형은 다음과 같다.

\chg{첫째, 실무 유의성을 나타내는 $f_1$은} \m{식~(\ref{eq:f1})에 나타낸 바와 같이} 평균 감소율 $\bar{r}(\kappa) = (1 - \bar{\eta}) \times 100$이
    목표 범위(5--30\%)에 있을 때 최대이며 15\%에서 정점을 갖는다.
    \begin{equation}
      f_1(\kappa) = \exp\!\left(-\frac{(\bar{r}(\kappa) - 15)^2}{2 \cdot 10^2}\right)
      \label{eq:f1}
    \end{equation}
    중심값 15\%는 RTO 감소율이 너무 작으면(5\% 미만) 도입 효과가 미미하고 너무 크면
    (30\% 초과) 비현실적이므로 목표 범위 5--30\%의 중앙값을 채택한 것이다. 분모의 $10^2$은
    가우시안 커널의 표준편차 $\sigma = 10$에 해당하며, 데이터에서 추정한 값이 아니라 목표 범위의
    구조에 맞추어 설정한 설계 파라미터이다. 목표 범위의 중심이 15\%이고 하한이 5\%이므로 중심에서
    하한까지의 거리는 $15 - 5 = 10$\%p이며, 이 거리를 $1\sigma$로 놓으면($\sigma = 10$) 가우시안
    커널의 확률적 해석과 자연스럽게 대응된다. $\sigma = 10$일 때 중심으로부터의 거리별 $f_1$ 값은
    \p{[}표~\ref{tab:f1-sigma10}\p{]}과 같다.
    \begin{table}[htbp]
    \centering
    \caption{$\sigma = 10$, 중심($\bar{r}=15\%$)에서 거리별 포화도 가중치 $f_1$}
    \label{tab:f1-sigma10}
    \small
    \begin{tabular}{c c c l}
    \toprule
    $\bar{r}$ & 중심 거리 & $f_1$ & 해석 \\
    \midrule
    15\% & $0\sigma$ & $e^{0} = 1.000$ & \m{최적점(정점)} \\
    5\% 또는 25\% & $1\sigma$ & $e^{-0.5} \approx 0.607$ & 목표 범위 내, 충분히 허용 \\
    35\% & $2\sigma$ & $e^{-2} \approx 0.135$ & 목표 초과, 강하게 억제 \\
    45\% & $3\sigma$ & $e^{-4.5} \approx 0.011$ & 사실상 배제 \\
    \bottomrule
    \end{tabular}
    \end{table}

\chg{둘째, 비포화를 나타내는 $f_2$는} \m{식~(\ref{eq:f2})에 나타낸 바와 같이} 포화율(시그모이드 출력 $> 0.95$인 비율)이 5\% 미만이면 1, 아니면 0이다.
    \begin{equation}
      f_2(\kappa) = \begin{cases}
        1.0 & \text{if sat}(\kappa) < 5\% \\
        0.0 & \text{otherwise}
      \end{cases}
      \label{eq:f2}
    \end{equation}
    시그모이드 출력이 0.95를 초과하면 입력 변화에 대한 출력 변화가 거의 없는 포화(saturation)
    영역에 진입하여 $\DRTO$가 입력 변수의 차이를 변별하지 못하므로, 포화 비율을 5\% 미만으로
    억제하는 이진 게이트로 설계하여 이 조건을 위반하는 $\kappa$를 후보에서 즉시 배제한다.

\chg{셋째, 효과 크기를 나타내는 $f_3$는} \m{식~(\ref{eq:f3})에 나타낸 바와 같이} C0 대비 Cohen's $d$의 \m{절댓값}이 $\geq 0.8$이면 1이다.
    \begin{equation}
      f_3(\kappa) = \begin{cases}
        1.0 & \text{if } |d(\kappa)| \geq 0.8 \\
        |d(\kappa)| / 0.8 & \text{otherwise}
      \end{cases}
      \label{eq:f3}
    \end{equation}
    식~(\ref{eq:f3})에서 C0 조건은 $\eta \equiv 1.000$의 결정론적 상수이므로 합동 표준편차 공식의 분모가
    인위적으로 축소되는 것을 피하기 위해 C1 자체의 표준편차를 분모로 사용하는 일표본 효과
    크기를 적용한다. $\bar{\eta}_{\text{C1}} = 0.853$, $s_{\text{C1}} = 0.125$이므로
    $d = (0.853 - 1.000)/0.125 = -1.176 \approx -1.180$이고, $|d| \geq 0.8$이므로 $f_3 = 1.0$이다. Cohen's $d$의 정의와 해석 기준은 4.8(조건간 효과 크기 분석)을 참조한다.

\chg{넷째, 설명력을 나타내는 $f_4$는} \m{식~(\ref{eq:f4})에 나타낸 바와 같이} $\eta_{C1}$에 대한 회귀 $R^2$이 $\geq 0.95$이면 1이다.
    \begin{equation}
      f_4(\kappa) = \begin{cases}
        1.0 & \text{if } R^2(\kappa) \geq 0.95 \\
        R^2(\kappa) / 0.95 & \text{otherwise}
      \end{cases}
      \label{eq:f4}
    \end{equation}

\p{[}표~\ref{tab:kappa-grid}\p{]}는 $\kappa \in \{0.2, 0.5, 1.0, 1.2, 1.5, 2.0, 3.0, 5.0\}$에
대한 격자 탐색(grid search) 결과를 정리한 것이다.
\chg{탐색은 표본 크기 $n = 10{,}000$, seed $= 42$에서 수행하였으며, 표의 $R^2_{(2\mathrm{q})}$는 2차 다항 회귀의 결정계수이다.}

\p{[}표~\ref{tab:kappa-grid}\p{]}에서 확인되듯이, $\kappa = 1.2$에서 종합 점수
$F(\kappa) = 1.000$으로 최대값을 달성한다.
이 지점에서 Cohen's $|d| = 1.180$(대효과), 포화율 $0.0\%$,
$R^2 = 0.999$, 단축률 $14.7\%$가 동시에 충족된다.

\begin{table}[htbp]
\centering
\caption{$\kappa$ 격자 탐색 \m{결과($n = 10{,}000$, seed $= 42$)}}
\label{tab:kappa-grid}
\small
\begin{tabularx}{\textwidth}{c c c c c c c}
\toprule
$\kappa$ & 평균 $\bar{\eta}_{\text{C1}}$ & 단축률 (\%) & $|d|_{\text{C1vsC0}}$
  & $R^2_{(2\mathrm{q})}$ & 포화율 (\%) & $F(\kappa)$ \\
\midrule
0.2  & 0.975 &  2.5 & 1.150 & 1.000 &  0.0 & \erf{0.813}{0.459} \\
0.5  & 0.937 &  6.3 & 1.150 & 1.000 &  0.0 & \erf{0.855}{0.684} \\
1.0  & 0.876 & 12.4 & 1.170 & 0.999 &  0.0 & \erf{0.957}{0.967} \\
\rowcolor{green!10}
1.2  & 0.853 & 14.7 & 1.180 & 0.999 &  0.0 & \textbf{1.000} \\
1.5  & 0.818 & 18.2 & 1.200 & 0.997 &  0.0 & \erf{0.947}{0.951} \\
2.0  & 0.765 & 23.5 & 1.230 & 0.993 &  0.0 & \erf{0.859}{0.697} \\
3.0  & 0.673 & 32.7 & 1.310 & 0.975 &  0.0 & \erf{0.858}{0.209} \\
5.0  & 0.538 & 46.2 & 1.490 & 0.908 &  3.9 & \erf{0.566}{0.007} \\
\bottomrule
\end{tabularx}
\end{table}


$\kappa < 1.2$인 영역에서는 효과 크기와 단축률이 부족하고,
$\kappa > 1.2$인 영역에서는 포화가 시작되어 $R^2$가 하락한다.
$\kappa = 1.2$는 이 두 영역의 균형점(trade-off point)에 정확히 위치한다.
\p{[}표~\ref{tab:kappa-optimal}\p{]}는 $\kappa^{*} = 1.2$에서 4개 부분 점수의 개별
달성값과 종합점수 $F^{*}$를 분해한 것이다.

\begin{table}[htbp]
\centering
\caption{$\kappa^{*} = 1.2$에서의 4가지 기준 달성값}
\label{tab:kappa-optimal}
\begin{tabular}{llrrl}
\toprule
\textbf{기준} & \textbf{지표} & \textbf{달성값} & \textbf{$f_i$ 점수} & \textbf{판정} \\
\midrule
$f_1$ (실무 유의성)  & 평균 감소율      & 14.74 & 0.9997 & 목표 15\%에 근접 \\
$f_2$ (비포화)       & 포화율           & 0.0   & 1.000 & 포화 완전 회피 \\
$f_3$ (효과 크기)    & Cohen's $d$     & $-1.18$ & 1.000 & 대효과 달성 \\
$f_4$ (설명력)       & $R^2$           & 0.999   & 1.000 & 고설명력 달성 \\
\midrule
\multicolumn{2}{l}{\textbf{종합점수 $F^*$}} & & \textbf{0.9997} & \textbf{유일 극대} \\
\bottomrule
\end{tabular}
\end{table}

\p{[}그림~\ref{fig:kappa-composite}\p{]}은 종합점수 $F(\kappa)$의 곡선을
도시한 것으로, $\kappa^{*} = 1.2$에서 최대값을 달성하는 것을 보여준다.
\chg{이 지점에서 효과 크기와 비포화, $R^2$, 단축률의 네 기준이 균형을 이루어 $F = 1.000$을 달성하며, $\kappa < 1.0$은 과소 반응 영역에, $\kappa > 2.0$은 조기 포화 영역에 해당한다.}

\begin{figure}[htbp]
\centering
\begin{tikzpicture}
\ifdefined\ERRATAFIX\def\FKymin{0}\def\FKcoords{(0.2, 0.459) (0.5, 0.684) (1.0, 0.967) (1.2, 1.000) (1.5, 0.951) (2.0, 0.697) (3.0, 0.209) (5.0, 0.007)}\else\def\FKymin{0.4}\def\FKcoords{(0.2, 0.813) (0.5, 0.855) (1.0, 0.957) (1.2, 1.000) (1.5, 0.947) (2.0, 0.859) (3.0, 0.858) (5.0, 0.566)}\fi
\begin{axis}[
  width=0.80\textwidth,
  height=0.48\textwidth,
  xlabel={$\kappa$},
  ylabel={$F(\kappa)$},
  xmin=0, xmax=5.5,
  ymin=\FKymin, ymax=1.05,
  grid=major,
  grid style={black},
  axis lines=left,
  every axis x label/.style={at={(axis description cs:1.02,0.0)},
    anchor=west, font=\small},
  every axis y label/.style={at={(axis description cs:-0.02,1.02)},
    anchor=south, font=\small},
  clip=false,
  xtick={0, 0.5, 1.0, 1.2, 1.5, 2.0, 3.0, 5.0},
  xticklabels={0, 0.5, 1.0, \textbf{1.2}, 1.5, 2.0, 3.0, 5.0}
]
  % 데이터 점 연결 곡선 (smooth)
  \addplot[thick, black, smooth, mark=*, mark size=2pt]
    coordinates {
      \FKcoords
    };

  % 최적점 강조
  \filldraw[red!80] (axis cs:1.2, 0.996) circle (3pt);
  \node[font=\small, text=black, anchor=south west] at (axis cs:1.3, 0.996)
    {$\kappa^{*} = 1.2$, $F = 1.000$};

  % 과소 반응 영역
  \draw[thick, blue!50, {Stealth[length=2mm]}-] (axis cs:0.1, 0.45) -- (axis cs:0.8, 0.45);
  \node[font=\scriptsize, text=black, anchor=north] at (axis cs:0.45, 0.45)
    {과소 반응};

  % 조기 포화 영역
  \draw[thick, orange!60, -{Stealth[length=2mm]}] (axis cs:2.5, 0.45) -- (axis cs:5.3, 0.45);
  \node[font=\scriptsize, text=black, anchor=north] at (axis cs:3.9, 0.45)
    {조기 포화};

\end{axis}
\end{tikzpicture}
\caption{$\kappa$ 격자 탐색에 의한 종합점수 $F(\kappa)$ 곡선}
\label{fig:kappa-composite}
\end{figure}
     % 2.9--2.10 수학적 성질·κ 최적화


%% ===================================================================
\section{연구가설}
\label{sec:hypotheses}
%% ===================================================================

$\DRTO$ 프레임워크의 이론적 타당성과 실증적 유효성을 검증하기 위해
\jrev{10개의 연구가설을 수립한다.} 가설은 시뮬레이션 검증\jrev{(H1--H3, 3건)},
회귀 검증(H4--H5, 2건), 강건성 검증(H6--H7, 2건),
이중채널 비선형 효과(H8, 1건), 전문가 평가(H9--H10, 2건)의 5개 범주로 구성된다.

\chg{각 가설은 서로 다른 이론적 근거에서 도출된다. H1에서 H3까지와 H8은 본 장에서 확립한 모델의 해석적 구조, 즉 시그모이드 바운딩의 부호 조건과 경계 보장 및 포화 성질에서 연역되고, H4와 H5는 입력 변수와 응답 사이의 비선형 관계를 완전 2차 다항식으로 근사하는 반응표면법~\citep{box1951experimental}에, H6과 H7은 몬테카를로 시뮬레이션의 독립 시드 설계와 계수 안정성 원리에, H9와 H10은 정량과 정성 방법을 결합하는 혼합연구방법의 삼각검증 원리~\citep{creswell2018designing}에 각각 근거를 둔다.}

\chg{각 가설에 대해서는 도출 배경과 핵심 명제, 검증 방향의 세 요소를 일관되게
제시한다. 도출 배경은 왜 이 가설이 필요한가를, 핵심 명제는 그 수학적이고 통계적인
정의를, 검증 방향은 어느 장에서 어떻게 검증되는가를 밝힌다. 시뮬레이션 검증에
해당하는 H1에서 H3까지는 모형의 기본 동작을, 회귀 검증인 H4와 H5는 변수 간 통계적
관계를, 강건성 검증인 H6과 H7은 결과의 재현성을, 이중채널의 H8은 본 연구의 핵심 신규
발견을, 전문가 평가인 H9와 H10은 외부 타당성을 각각 담당한다.
\secref{sec:ch4_hypothesis_summary}에서 10개 가설이 모두 채택(PASS)으로 종합되며, 그
정량 근거는 가설별 정의와 시뮬레이션 및 회귀 결과로 뒷받침된다.}


%% -------------------------------------------------------------------
%% -------------------------------------------------------------------

\chg{가설 H1(관측성 효과)은 관측성이 활성화된 조건(C1)에서 $\DRTO$가 정적 기준선(C0)보다 작거나 같다는 것이다.}
\[
  H_1:\;\; \DRTO_{\textbf{C1}} \leq \DRTO_{\textbf{C0}} = R_0
  \quad (\text{준수율 } 100\%)
\]

H1은 관측성의 효과 방향성을 검증한다.
수학적으로, C1 조건에서 $z_O \geq 0$이므로 $S(z_O) \geq 0$이 되어
$E_{O,\text{bnd}} = -R_0 \cdot S(z_O) \leq 0$이 항상 성립한다.
따라서 $\DRTO_{\text{C1}} = R_0 + E_{O,\text{bnd}} \leq R_0$가
모델 구조에 의해 해석적으로 보장된다.


\chg{가설 H2(시그모이드 경계 보장)는 모든 조건(C0에서 C3까지)에서 $\DRTO$가 물리적 경계 내에 위치한다는 것이다.}
\[
  H_2:\;\; 0 < \DRTO < R_{\max} = 4R_0 = 24.0\text{h}
  \quad (\text{준수율 } 100\%)
\]

H2는 개별 시그모이드 바운딩 모델의 핵심 설계 목표인
물리적 경계 보장을 검증한다. (명제 2.1)에 의해
이 경계는 $\kappa$에 무관하게 항상 성립한다.



\chg{가설 H3(조건 순서 보존)는 네 조건이 다음의 세 가지의 순서 관계를 유지한다는 것이다. $H_{3\text{-}a}$는
$\DRTO_{\textbf{C1}} \leq \DRTO_{\textbf{C0}}$(준수율 $100\%$), $H_{3\text{-}b}$는
$\DRTO_{\textbf{C1}} \leq \DRTO_{\textbf{C2}}$(준수율 $\geq 99\%$), $H_{3\text{-}c}$는
$\DRTO_{\textbf{C2}} \leq \DRTO_{\textbf{C3}}$(조건부, 백분위 $\geq$ P85.8)을 의미한다.}

H3는 조건 간 $\DRTO$의 크기 순서가 이론적 예측과 일치하는지를 검증한다.
H3-a는 H1과 동치이며 모델 구조에 의해 보장된다.
H3-b는 불확실성의 도입이 관측성의 하향 효과를 상쇄하고 초과함을 의미하며,
$N(3,1)$의 음수 표본(약 $0.135\%$)에서 $E_{U,\text{bnd}} < 0$이 되는 경우를
허용하기 위해 준수율 기준을 100\%가 아닌 99\% 이상으로 설정한다.
H3-c는 파레토 분포 특성이 가우시안 대비 높은 $\DRTO$를 산출함을
검증하나, 이 효과는 P85.8 CDF 교차점 이상의 극단 영역에서만 발현되므로
조건부 가설로 설정한다.


%% -------------------------------------------------------------------
%% -------------------------------------------------------------------

\chg{가설 H4(회귀 설명력)는 교호 작용을 포함한 2차 회귀 모델의 결정계수가 모든 조건에서 $R^2_{(2\mathrm{q})} \geq 0.95$를 달성한다는 것이다.}
\[
  \jrev{H_4}:\;\; R^2_{(2\mathrm{q})} \geq 0.95
  \quad \forall\; \textbf{C1}, \textbf{C2}, \textbf{C3}
\]

\erf{H5}{H4}는 $\DRTO$ 모델의 비선형 응답 곡면이 2차 다항식으로 충분히
근사될 수 있는지를 검증한다. 2차 회귀는 1차항($k_O$, $O$, $U$),
2차항($k_O^2$, $O^2$, $U^2$), 교호작용항($k_O \cdot O$, $k_O \cdot U$,
$O \cdot U$)을 포함하는 완전 2차 다항 모델이다.
\chg{이러한 모델 선택의 이론적 근거는 \citet{box1951experimental}이 정립한 반응표면법에 있다. 반응표면법은 입력 변수와 응답 사이의 미지의 비선형 관계를 국소적으로 곡률까지 포착하는 완전 2차 다항식으로 근사하는 접근이 타당함을 보였으며, 본 연구의 응답 곡면 또한 시그모이드 바운딩에서 비롯한 매끄러운 곡률 구조를 가지므로 이 근사가 적합하다.}


\chg{가설 H5(교호 작용 유의성)는 교호작용항(interaction terms)이 회귀 모델의 설명력에 유의하게 기여한다는 것이다.}
\[
  \jrev{H_5}:\;\; \Delta R^2_{(\text{교호작용항 추가})} > 0,\quad F\text{-검정에서 } p < 0.05
\]

H5는 관측성과 불확실성이 독립적으로 작용하는 것이 아니라,
두 요소의 곱(교호 작용)이 $\DRTO$에 추가적인 설명력을 제공함을 검증한다.
특히 $k_O \cdot O_{\text{raw}}$ 교호작용항은 C1에서 $R^2$를 결정적으로 향상시킨다.


%% -------------------------------------------------------------------
%% -------------------------------------------------------------------

\chg{가설 H6(비중첩 대역 시드 독립성)은 비중첩 대역(non-overlapping band) 시드에서 파생된 입력 변수 $k_O$, $O_{\text{raw}}$, $U_{\text{raw}}$ 간 상관계수의 \m{절댓값}이 0.02 미만이라는 것이다.}
\[
  \jrev{H_6}:\;\; |r(X_i, X_j)| < 0.02 \quad \forall\; i \neq j
\]

H6은 시뮬레이션의 입력 변수들이 통계적으로 독립임을 검증한다.
마스터 시드 42로부터 비중첩 대역 규칙에 따라 파생된 시드 간
상관계수가 무시 가능한 수준임을 확인한다.


\chg{가설 H7(다중 시드 강건성)은 다중 시드($s = 0, 1, \ldots, 9{,}999$) 시뮬레이션에서 회귀 계수가 안정적이라는 것이다.}
\[
  \jrev{H_7}:\;\; \mathrm{CV}(\beta_{k_O}) < 0.05
\]

H7은 seed = 42에서 도출된 결과가 특정 시드에 의한 우연적 산물이 아님을
검증한다. 10{,}000개의 독립 시드 각각에서 $n = 10{,}000$ 시뮬레이션을
수행하여 총 $10^8$회의 $\DRTO$ 계산을 실시하고,
회귀계수의 분포를 산출한다.


%% -------------------------------------------------------------------
%% -------------------------------------------------------------------

\chg{가설 H8(이중채널 비선형 감쇠)은 C3 꼬리 영역($\geq \text{P85.8}$)에서 관측성 가중치 $k_O$의 회귀 계수 크기가 핵심 영역($< \text{P85.8}$) 대비 감소한다(감쇠비 $< 1$)는 것이다.}
\[
  \jrev{H_8}:\;\; \frac{|\beta_{k_O}^{\text{Tail}}|}{|\beta_{k_O}^{\text{Core}}|} < 1
\]

H8은 본 연구에서 새롭게 발견한 이중 채널 비선형 상호작용(dual-channel nonlinear
interaction) 메커니즘을 형식화한 가설이다.
``이중 채널''이란 관측성 채널($z_O$)과 불확실성 채널($z_U$)의 두 경로를 의미하며,
두 채널이 시그모이드 함수의 비선형성을 통해 결합될 때
극단 영역에서 시그모이드 포화(saturation)로 인해 관측성의 한계 효과가
오히려 감소하는 현상을 지칭한다.
감쇠비의 실증값과 P85.8 분할 회귀 분석 상세 내용은 4.10(이중채널 감쇠 메커니즘)에서 다루며,
이는 극단적 불확실성 하에서 관측성 투자의 한계 효과가 체감함을 의미한다.
이러한 비선형 감쇠 메커니즘은 조직의 위험 관리 전략에 대한
중요한 정책적 시사점을 제공하며 5.3(실무적 기여)에서 상세히 논의한다.


%% -------------------------------------------------------------------
%% -------------------------------------------------------------------

\chg{아홉째 가설 H9(전문가 정량적 검증)는 BCM/DR 전문가 설문조사에서 전체 영역 평균이 5점 리커트 척도 기준 4.0 이상이라는 것이다.}
\[
  \jrev{H_{9}}:\;\; \bar{X}_{\text{expert}} \geq 4.0 \;/\; 5.0
\]


\chg{열째 가설 H10(전문가 정성적 합의)은 전문가 심층 인터뷰에서 도출된 핵심 주제에 대한 동의율이 $80\%$ 이상이라는 것이다. 각 주제 $j$에 대한 동의 여부를 이진 변수 $a_j \in \{0, 1\}$로 정의하면 $a_j = 1$은 동의를, $a_j = 0$은 비동의를 가리킨다.}
\[
  \jrev{H_{10}}:\;\; A_{\text{agree}} \geq 0.80, \quad
  A_{\text{agree}} = \frac{1}{N_{\text{topic}}}\sum_{j=1}^{N_{\text{topic}}} a_j
\]
\chg{여기서 $A_{\text{agree}}$는 전체 $N_{\text{topic}}$개 주제 중 전문가가 동의한 주제의 비율을 의미한다.}

H9와 H10은 $\DRTO$ 모델의 실무적 타당성을 현장 전문가의 판단을 통해
삼각검증(triangulation)하기 위한 가설로서, \chg{정량적 방법과 정성적 방법을 결합하여 단일 방법의 한계를 보완하는 혼합연구방법의 삼각검증 원리~\citep{creswell2018designing}에 이론적 근거를 둔다.}
정량적 설문(H9)은 4개 평가 영역, \chg{곧 모델 타당성과 파라미터 적정성,
실무 적용성, 산업적용 및 $\eta$ 활용에 대한} 27개 문항으로 구성되며,
정성적 인터뷰(H10)는 반구조화 방식으로 수행한다.
\chg{H9와 H10의 검증 결과 상세내용은 4.11(전문가 패널 검증)에서 다룬다.}


%% -------------------------------------------------------------------
%% -------------------------------------------------------------------

\p{[}표~\ref{tab:hypotheses-summary}\p{]}에 \jrev{10개} 가설의 정의, 판정 기준, 관련 조건을 종합하였다.
\chg{[표 2-16]에서 H7의 $10^8$은 다중 시드 시뮬레이션의 총 횟수를 가리킨다.}

\begin{table}[htbp]
\centering
\caption{\jrev{10개 연구가설 종합}}
\label{tab:hypotheses-summary}
\small
\begin{tabularx}{\textwidth}{c c X l c}
\toprule
\textbf{가설} & \textbf{범주} & \textbf{핵심 내용} & \textbf{판정 기준} & \textbf{조건} \\
\midrule
H1 & 시뮬레이션 & $\DRTO_{\text{C1}} \leq R_0$ (관측성 효과) & 준수율 100\% & C0, C1 \\
H2 & 시뮬레이션 & $0 < \DRTO < R_{\max} = 24.0$h (경계 보장) & 준수율 100\% & C0--C3 \\
H3 & 시뮬레이션 & C0--C3 순서 보존 (a/b/c) & 준수율 기준 & C0--C3 \\
\midrule
H4 & 회귀 & $R^2_{(2\mathrm{q})} \geq 0.95$ & C1/C2/C3 전체 & C1--C3 \\
H5 & 회귀 & 교호작용항 \m{유의성($\Delta R^2$)} & 기여도 검정 & C1--C3 \\
\midrule
H6 & 강건성 & 시드 간 $|r| < 0.02$ (독립성) & 상관계수 검정 & 전체 \\
H7 & 강건성 & 다중 시드 계수 \m{안정성($10^8$회)} & 계수 CV 기준 & 전체 \\
\midrule
H8 & 이중채널 & Tail/Core $k_O$ 감쇠비 $< 1$ & P85.8 분할 회귀 & C3 \\
\midrule
H9 & 전문가 & 설문 평균 $\geq 4.0/5.0$ & $t$-검정 & 전체 \\
H10 & 전문가 & $A_{\text{agree}} \geq 0.80$ (15개 하위 주제) & 합의 기준 & 전체 \\
\bottomrule
\end{tabularx}
\end{table}
          % 2.11 \jrev{10개 연구가설}
